\begin{textsample}{POS Dim 4 – 2014\_c6 – Score 3.00 – 2014\_c6/enade\_2014\_fis\_e35869ce3f.txt}  \label{ex:f4_pos_079}
CAPÍTULO 3 ANÁLISE TÉCNICA DA PROVA Este capítulo tem por objetivo apresentar o desempenho dos estudantes concluintes de Física ( Bacharelado e Licenciatura ) no ENADE/2014.
Para isso, foram calculadas as estatísticas básicas da prova em seu todo ( seção 3.1.1 ), bem como as estatísticas dos componentes relacionadas ao Componente de Formação Geral ( seção 3.1.2 ) e ao de Conhecimento Específico da Área ( seção 3.1.3 ).
Dadas as suas características, foram analisadas, em separado, as questões objetivas ( seção 3.2 ) e as questões discursivas ( seção 3.3 ) para concluintes de Bacharelado e de Licenciatura.
Para as questões objetivas, foram disponibilizados os índices de facilidade e o de discriminação bisserial, também em separado, para os Componentes de Formação Geral ( seção 3.2.1 ) e de Conhecimento Específico ( seção 3.2.2 ), sendo a análise da segunda seção separada por concluintes de Bacharelado e de Licenciatura.
De cada componente, uma das questões foi escolhida para exemplificar a análise gráfica, relacionando as alternativas escolhidas pelos estudantes ( inclusive o gabarito ) com o número de acertos no componente aferido, como já dito anteriormente, agregados por Bacharelado e Licenciatura.
O Anexo I apresenta a íntegra da análise gráfica para todas as questões objetivas agregadas e também separadas por Bacharelado e Licenciatura, a partir da questão 26.
Para cada uma das questões discursivas, os conteúdos dos tipos mais comuns de respostas dos estudantes são apresentados e comparados com o padrão de respostas esperado ( ver Anexo VIII com o padrão de respostas ).
Tomando como base as duas questões discursivas do Componente de Formação Geral, a seção 3.3.1.6 apresenta comentários sobre a correção das respostas com respeito à Língua Portuguesa.
Na Tabela 3.1 são apresentados o tamanho da população inscrita e de presentes, e as seguintes estatísticas das notas19 : média do desempenho na prova, erro padrão da média, desvio padrão, nota mínima, mediana e nota máxima.
Tais estatísticas contemplam o total de estudantes concluintes da Área de Física ( Bacharelado e Licenciatura ) inscritos e presentes à prova do ENADE/2014, tendo em vista as agregações por Grandes Regiões e o país como um todo.
As estatísticas das demais tabelas ímpares serão baseadas na mesma população da Tabela 3.1 e, portanto, do mesmo tamanho.
As tabelas pares são desagregadas por Categoria Administrativa e Organização Acadêmica da IES.
As estatísticas dessas tabelas pares serão baseadas na mesma população da Tabela 3.2. 19 Uma definição dessas estatísticas pode ser encontrada no Glossário.
Em relação aos gráficos de distribuição de notas, o intervalo considerado foi de 10 unidades, aberto à esquerda e fechado à direita, com exceção do primeiro intervalo, [ 0 ; 10 ], fechado em ambos os extremos.
Para os gráficos de distribuição das notas das questões discursivas, foram consideradas mais duas categorias : questão em branco20 e nota zero.
Todos os gráficos de distribuição de notas permitem a comparação dos resultados por habilitação : Bacharelado e Licenciatura. 3.1 ESTATÍSTICAS BÁSICAS DA PROVA Esta seção apresenta estatísticas selecionadas e histogramas da nota geral ( 3.1.1 ) e de cada componente : Formação Geral ( 3.1.2 ) e Conhecimento Específico ( 3.1.3 ).
São também apresentadas as estatísticas selecionadas de subpopulações caracterizadas por Grande Região, Categoria Administrativa e Organização Acadêmica. 3.1.1 Estatísticas Básicas Gerais A Tabela 3.1 apresenta as estatísticas básicas da prova, por grande Região, dos estudantes concluintes de Física ( Bacharelado e Licenciatura ).
A população total de inscritos foi de 4.200.
Destes, 3.280 estiveram presentes, sendo 21,9 \% o índice de não comparecimento.
A Região de maior abstenção foi a Norte ( 31,1 \% ), e a de menor abstenção foi a Sul ( 14,1 \% ).
A média das notas da prova como um todo ( nas seções seguintes serão analisados os Componentes de Formação Geral e de Conhecimento Específico ) foi 39,6, sendo que os estudantes da região Norte obtiveram a média mais baixa ( 34,5 ), e os da região Sul obtiveram a média mais alta ( 42,9 ).
As demais médias foram : 37,7 na região Nordeste ; 42,5 na região Sudeste ; e 40,8 na região Centro-Oeste.
O desvio padrão para o Brasil como um todo foi 14,1, sendo o maior desvio padrão encontrado na região Sudeste ( 14,9 ) e o menor, na região Norte ( 10,9 ), indicando uma dispersão um pouco menor das notas desta última região.
A região que obteve a maior nota máxima foi a Sul ( 86,4 ), ao passo que a região que atingiu a menor nota máxima foi a Centro-Oeste ( 72,7 ).
A mediana do Brasil como um todo foi 38,6, sendo a maior mediana obtida na região Sul ( 42,5 ), e a menor obtida na Norte ( 33,7 ).
A nota mínima foi zero em todas as regiões, exceto na região Norte ( 8,5 ).
Considerando -se as notas segundo Grande Região, observa -se que existe diferença estatisticamente significativa ao nível de 95 \% entre a menor média, obtidas na região Norte ( 34,5 ), e as médias de todas as demais regiões.
A menor obtida na Norte ( 33,7 ).
A nota mínima foi zero em todas as regiões, exceto na região Norte ( 8,5 ).
Considerando -se as notas segundo Grande Região, observa -se que existe diferença estatisticamente significativa ao nível de 95 \% entre a menor média, obtidas na região Norte ( 34,5 ), e as médias de todas as demais regiões.
O comportamento das notas dos estudantes de todo o Brasil pode ser observado no Gráfico 3.1 que apresenta um histograma com a distribuição das mesmas em colunas diferentes para estudantes de Bacharelado e de Licenciatura.
As distribuições são unimodais.
O intervalo modal dos concluintes de Bacharelado e de Licenciatura é o mesmo, ( 30 ; 40 ].
A Tabela 3.2 apresenta as estatísticas básicas da prova, desagregadas por Categoria Administrativa e por Organização Acadêmica.
Da população total de inscritos, 494 são de IES Privadas e 3.706 de IES Públicas.
Em relação à Organização Acadêmica, a maior participação foi obtida por estudantes de Universidades ( 4.038 ), seguido de Faculdades ( 85 ) e Centros Universitários ( 77 ).
A Categoria Administrativa de maior abstenção foi a Privada ( 23,5 \% ), e entre as Organizações Acadêmicas foi a dos Centros Universitários ( 23,4 \% ), ambos acima da média nacional de 21,9 \%.
A média das notas da prova como um todo foi 39,6.
Em relação à Categoria Administrativa, os estudantes das IES Públicas obtiveram média mais alta ( 39,7 ), e os das IES Privadas obtiveram média mais baixa ( 39,0 ), que a média nacional.
Observa -se que não existe diferença estatisticamente significativa entre as médias das notas das IES Públicas e Privadas.
A diferença entre as médias das regiões Sul e Norte ( 8,4 ), a maior e a menor média, é superior à diferença entre IES Públicas e Privadas ( 0,7 ), caracterizando uma maior diversidade regional do que administrativa.
No tocante à Organização Acadêmica, apenas as Faculdades obtiveram média mais baixa ( 35,2 ) que a nacional.
As demais médias foram superiores à média nacional : 43,5 nos Centros Universitários e 39,7 nas Universidades.
O desvio padrão para as IES Públicas ( 14,2 ) e para os Centros Universitários ( 18,0 ) foi superior ao do Brasil como um todo ( 14,1 ), indicando uma dispersão um pouco maior das notas nesta Categoria Administrativa e nesta Organização Acadêmica.
Constata -se que existe diferença estatisticamente significativa ao nível de 95 \% nas médias das notas dos estudantes provenientes de Faculdades em relação aos de Centros Universitários e Universidades. 3.1.2 Estatísticas Básicas no Componente de Formação Geral A Tabela 3.3 apresenta as estatísticas básicas em relação ao componente da prova que avalia a Formação Geral dos estudantes concluintes.
Os estudantes de todo Brasil obtiveram desempenho médio de 55,7.
Quanto à variabilidade, o desvio padrão das notas dos estudantes do Brasil, como um todo, foi 19,1.
A maior média foi obtida na região Sul ( 58,4 ), e a menor, na região Norte ( 52,8 ).
As demais médias foram : 54,5 na região Nordeste ; 57,0 na região Sudeste ; e 56,8 na região Centro-Oeste.
Já o maior desvio padrão foi obtido na região Sudeste ( 19,6 ), e o menor, na região Norte ( 17,3 ).
Os demais desvios padrões foram : 19,0 na região Nordeste, 19,3 na região Sul e 19,5 na região Centro-Oeste.
A maior nota no Componente de Formação Geral da prova do ENADE/2014 foi 98,4, obtida por, pelo menos, um estudante nas regiões Sudeste e Centro-Oeste.
A menor nota máxima foi obtida na região Norte ( 93,6 ).
A mediana do Brasil, como um todo, foi 57,8, sendo a menor mediana encontrada na região Norte ( 53,6 ) ; e a maior encontrada na região Sul ( 61,6 ).
A nota mínima nesta parte foi zero em todas as regiões, exceto na região Norte ( 7,5 ).
Considerando -se as notas segundo Grande Região, observa -se que existe diferença estatisticamente significativa entre a maior média das notas do Componente de Formação Geral, obtida na região Sul ( 58,4 ), e na região Norte ( 52,8 ), a menor média. encontrada na região Sul ( 61,6 ).
A nota mínima nesta parte foi zero em todas as regiões, exceto na região Norte ( 7,5 ).
Considerando -se as notas segundo Grande Região, observa -se que existe diferença estatisticamente significativa entre a maior média das notas do Componente de Formação Geral, obtida na região Sul ( 58,4 ), e na região Norte ( 52,8 ), a menor média.
O Gráfico 3.2 propicia a avaliação do desempenho dos estudantes no Componente de Formação Geral a partir do histograma da distribuição das notas correspondentes ao Bacharelado e à Licenciatura.
A distribuição para o Bacharelado é bimodal, com a modas, respectivamente, nos intervalos ( 50 ; 60 ] e ( 70 ; 80 ].
É unimodal, com a moda no intervalo ( 60 ; 70 ], para a Licenciatura.
Nota -se, ainda, que no Gráfico 3.2 as notas apresentam uma maior dispersão do que no Gráfico 3.1 ( distribuição das notas da prova ), confirmado pela comparação dos desvios padrões : 14,1 para a nota da prova como um todo e 19,1 para o Componente de Formação Geral.
Na Tabela 3.4 são apresentadas as informações referentes ao desempenho dos concluintes do Componente de Formação Geral, em diferentes agregações : Categoria Administrativa e Organização Acadêmica.
Observa -se que a maior média foi obtida por estudantes de IES Privadas ( 56,3 ), com uma diferença que não é estatisticamente significativa da obtida por estudantes de IES Públicas ( 55,6 ).
Considerando -se o tipo de Organização Acadêmica, nota -se que não há diferença estatisticamente significativa entre as médias dos estudantes de Centros Universitários ( 58,1 ), e de Universidades ( 55,7 ) e de Faculdades ( 53,7 ). 3.1.3 Estatísticas Básicas do Componente de Conhecimento Específico A Tabela 3.5 apresenta as estatísticas básicas referentes ao Componente de Conhecimento Específico da área de Física ( Bacharelado e Licenciatura ).
A média do desempenho dos estudantes do Brasil, como um todo, foi 34,3 para este componente.
A maior média foi obtida na região Sul ( 37,7 ), e a menor, na região Norte ( 28,4 ).
As demais médias foram : 32,1 na região Nordeste ; 37,6 na região Sudeste ; e 35,5 na região Centro-Oeste.
Quanto à variabilidade das notas, o desvio padrão do Brasil, como um todo, foi 15,2, sendo o maior desvio padrão observado na região Sudeste ( 16,1 ), e o menor, na região Norte ( 11,8 ).
Os demais desvios foram : 14,4 na região Nordeste, 15,5 na região Sul e 15,7 na região Centro-Oeste.
A mediana das notas dos estudantes de todo o Brasil foi 32,3.
A maior mediana ocorreu na região Sul ( 36,0 ), e a menor, na região Norte ( 26,9 ).
As demais medianas foram : 31,1 na região Nordeste ; 35,8 na região Sudeste ; e 35,0 na região Centro-Oeste.
A nota máxima do Brasil, como um todo, foi 89,3, sendo obtida por, pelo menos, um estudante na região Sul.
As demais notas máximas foram : 79,1 na região Norte ; 83,6 na região Nordeste ; 87,8 na região Sudeste ; e 81,1 na região Centro-Oeste.
A nota mínima foi zero em todas as regiões, exceto na região Norte ( 2,5 ).
Observa -se que existe diferença estatisticamente significativa da menor média das notas do Componente de Conhecimento Específico da região Norte ( 28,4 ) em relação às demais regiões.
O intervalo de confiança para a média da região Centro-Oeste ( 4,6 ) é maior do que de todas as outras.
Assim como os Gráficos 3.1 e 3.2, o Gráfico 3.3, apresentado a seguir, permite uma avaliação do desempenho de concluintes dos cursos de Bacharelado e de Licenciatura em relação ao Componente de Conhecimento Específico com um histograma da distribuição das notas correspondentes.
Dentre as três distribuições apresentadas, esta é a mais concentrada nas notas baixas.
Ambas são distribuições unimodais.
O grupo modal é o ( 30 ; 40 ] para o Bacharelado e o ( 20 ; 30 ] para a Licenciatura.
A Tabela 3.6 apresenta uma comparação dos resultados em relação à Categoria Administrativa e à Organização Acadêmica, agora levando em conta o desempenho de estudantes do Componente de Conhecimento Específico da prova.
No que se refere à Organização Acadêmica, a maior média foi dos Centros Universitários ( 38,6 ), vindo a seguir a das Universidades ( 34,3 ) e, depois, a das Faculdades ( 29,0 ).
O maior desvio padrão, e acima do valor para o Brasil, como um todo, foi o dos Centros Universitários ( 19,8 ).
As Universidades obtiveram a maior nota máxima ( 89,3 ) seguidas dos Centros Universitários ( 87,8 ) e das Faculdades ( 51,2 ).
Os Centros Universitários obtiveram a maior mediana ( 33,5 ) seguidos das Universidades ( 32,3 ) e das Faculdades ( 30,4 ).
A nota mínima foi zero para as Universidades e 8,9 nas demais Organizações Acadêmicas.
Quanto à Categoria Administrativa, observa -se um comportamento semelhante àquele da parte de Formação Geral e à prova como um todo, ou seja, não existe diferença estatisticamente significativa entre as médias das IES Públicas ( 34,4 ) e IES Privadas ( 33,2 ).
Neste caso, a maior média foi obtida por estudantes de IES Públicas de ensino.
Observa -se que existem diferenças estatisticamente significativas ao nível de 95 \% no Componente de Conhecimento Específico entre a nota média das Faculdades ( 29,0 ) em relação aos Centros Universitários ( 38,6 ) e Universidades ( 34,4 ). 3.2 ANÁLISE DAS QUESTÕES OBJETIVAS Esta seção apresenta estatísticas para o conjunto de questões objetivas e uma análise gráfica de questões selecionadas dos componentes de Formação Geral ( 3.2.1 ), de Conhecimento Específico Bacharelado ( 3.2.2 ) e de Conhecimento Específico Licenciatura ( 3.2.3 ).
São também apresentadas e comparadas as médias das subpopulações caracterizadas por Grande Região. 3.2.1 Componente de Formação Geral A Tabela 3.7 apresenta as estatísticas básicas relativas às oito questões objetivas do componente da prova que abrange a Formação Geral dos estudantes.
A média do Brasil foi 58,9.
A menor média foi encontrada na região Norte ( 55,5 ), e a maior, na região Sul ( 62,2 ).
As demais médias foram : 57,3 na região Nordeste ; 60,7 na região Sudeste ; e 60,5 na região Centro-Oeste.
O desvio padrão do Brasil foi 22,3, sendo o maior desvio padrão encontrado na região Centro-Oeste ( 24,5 ), e o menor, na região Norte ( 20,6 ).
Os demais desvios foram : 21,9 na região Nordeste ; 22,6 na região Sudeste ; e 22,3 na região Sul.
As medianas ( 62,5 ), as notas máximas ( 100,0 ) e as notas mínimas ( 0,0 ) foram iguais para todas as regiões, sem exceção.
A Tabela 3.8 apresenta o índice de facilidade e o índice de discriminação ( ponto bisserial ) para cada uma das questões objetivas do Componente de Formação Geral.
Quanto ao índice de facilidade, foram usadas as seguintes cores para diferenciar o nível de dificuldade da questão : • Azul para as questões classificadas com índice muito fácil ( \&gt ; =0,86 ), verde para as questões classificadas com índice fácil ( 0,61 a 0,85 ), amarelo para as questões classificadas com médio ( 0,41 a 0,60 ), vermelho para as questões classificadas com difícil ( 0,16 a 0,40 ) e roxo para as questões classificadas com muito difícil ( \&lt ; =0,15 ).
Já quanto ao índice de discriminação, foram usadas as seguintes cores para qualificar a questão :  As questões classificadas com índice fraco receberam a cor vermelho ( \&lt ; =0,19 ), as classificadas com médio receberam a cor amarelo ( 0,20 a 0,29 ), as classificadas com bom receberam a cor verde ( 0,30 a 0,39 ) e as classificadas com muito bom ( \&gt ; =0,40 ) receberam a cor azul.
As questões objetivas do Componente de Formação Geral, segundo o índice de facilidade, foram assim avaliadas : das oito questões, nenhuma teve o índice de facilidade classificado como muito fácil.
Quatro questões foram tidas como fácil, por terem índice de acertos situado na faixa entre 0,61 e 0,85 ( de 61,0 \% a 85,0 \% de acertos ).
Três questões foram consideradas de dificuldade médio, situando -se no intervalo entre 0,41 e 0,60 do índice de facilidade, ou seja, houve entre 41,0 \% e 60,0 \% de acertos, enquanto uma questão foi classificada na categoria difícil, situando -se no intervalo entre 0,16 e 0,40.
Por fim, nenhuma das questões apresentou menos de 15 \% de acertos, razão pela qual seria classificada como muito difícil.
Como já comentado, para análise das questões objetivas, relativas à Formação Geral, segundo o poder de discriminação, utilizou -se o índice de discriminação ( ponto bisserial ).
Nesta análise, as questões foram assim avaliadas : sete das oito questões apresentaram índice acima ou igual a 0,40 e, assim, foram classificadas com índice muito bom para esse grupo de estudantes.
Uma questão teve índice de discriminação bom, com valor entre 0,30 e 0,39, para esse grupo de estudantes.
Nenhuma questão teve nível médio ou fraco de discriminação para esse grupo de estudantes.
O índice de facilidade variou de 0,37 a 0,80, e o de discriminação, de 0,36 a 0,56.
As sete questões com índice de discriminação muito bom figuraram entre os diversos níveis de dificuldade desse conjunto : quatro classificadas na categoria fácil ( questões 3, 4, 6 e 8 ) do índice de facilidade, duas, na categoria médio ( questões 1 e 7 ) e uma, na categoria difícil ( questão 5 ).
Em particular, a questão 7 foi a que apresentou o maior poder discriminatório, com índice 0,56, porém foi considerada média em termos de facilidade, com uma proporção de 0,56 acertos.
O máximo de acertos foi alcançado pela questão 4 com um índice de facilidade de 0,80.
A questão de número 1 apresentou um índice de facilidade de 0,51, ou seja, pouco mais da metade dos estudantes conseguiu resolvê -la, dentro do universo de participantes.
Seu índice de discriminação foi muito bom ( 0,41 ).
Já a questão 5 obteve índice de discriminação muito bom, 0,44, porém seu índice de facilidade foi difícil ( 0,37 ).
O Gráfico 3.4, para exemplificar, analisa o comportamento da questão de número 7 de Formação Geral.
Trata -se de uma questão considerada mediana, com relação a facilidade e a que obteve o maior índice de discriminação dessa parte da prova.
Neste gráfico, cada uma das cinco curvas representa o percentual de respostas em determinada alternativa da questão, em função do número de acertos dos estudantes nessa parte da prova ( Formação Geral/Múltipla Escolha ), antes de possíveis eliminações pelo critério do ponto bisserial.
A curva em vermelho corresponde à alternativa E, a correta para esta questão.
Assim, observa -se que entre os estudantes com menor número de acertos, nessa parte do exame, a situação mais frequente foi a escolha de uma das alternativas incorretas : a alternativa C ( em preto ) ou D ( em roxo ).
À medida em que o número de acertos aumenta, indicando desempenho melhor nessa parte da prova, aumenta concomitantemente a proporção de estudantes que selecionaram a alternativa correta E, atingindo 100 \% para os estudantes com 8 acertos.
Essa análise permite verificar como a questão discriminou os grupos de desempenho, justificando o alto índice obtido na questão.
Os gráficos relativos às demais questões de Formação Geral constam do Anexo I. 3.2.2 Componente de Conhecimento Específico – Física ( Bacharelado ) A Tabela 3.9 apresenta as estatísticas básicas em relação às questões objetivas do Componente de Conhecimento Específico da prova de Física ( Bacharelado ), por Grande Região.
A média do Brasil deste componente foi de 32,2.
A menor média foi observada na região Norte ( 29,9 ), e a maior, na região Centro-Oeste ( 38,1 ).
O desvio padrão de todo o Brasil foi 14,4, sendo o menor desvio padrão encontrado na região Norte ( 9,8 ), e o maior, na região Centro-Oeste ( 15,8 ).
A mediana de todo o Brasil foi 31,8, a mesma encontrada nas regiões Nordeste, Sudeste e Sul.
A região Norte ( 29,6 ) apresentou valor menor para a mediana, enquanto a região Centro-Oeste ( 36,4 ), valor maior.
A nota máxima da prova foi 77,3, obtida nas questões objetivas do Componente de Conhecimento Específico, por pelo menos, um estudante na região Sudeste.
Nas demais regiões, a nota máxima da prova foi : 50,0 na região Norte ; 72,7 na região Nordeste ; e 68,2 nas regiões Sul e Centro-Oeste.
Em quase todas as regiões a nota mínima foi zero, exceto nas regiões Norte ( 13,6 ) e Centro-Oeste ( 9,1 ).
A Tabela 3.10 apresenta os índices de facilidade e discriminação ( ponto bisserial ) das questões objetivas do Componente de Conhecimento Específico para os estudantes de Física ( Bacharelado ).
Para facilitar a diferenciação das questões usaram -se as mesmas cores da Tabela 3.8 para as diferentes classificações dos índices de facilidade e de discriminação.
Dentre as questões objetivas da parte da prova relativa ao Componente de Conhecimento Específico, uma questão foi anulada pela Comissão, a de número 33.
Desse modo, a classificação, quanto ao índice de facilidade, foi estabelecida com base em 26 questões.
A partir dos índices obtidos, pode -se concluir que um pouco mais de três quartos das questões objetivas da prova foi considerado, pelo menos, difícil : das 26 questões válidas, 20 foram classificadas como difícil ( a classificação modal ) ou como muito difícil.
Uma questão foi classificada como fácil, e outras cinco consideradas como médio.
Já quanto aos índices de discriminação das questões objetivas do Componente de Conhecimento Específico da prova, tem -se como resultado a seguinte classificação : nove das 26 questões foram consideradas como boas, enquanto quatro delas tiveram índice de discriminação muito bom.
Assim, para a metade das questões – 13 em 26 – os índices de discriminação foram bom ou muito bom.
Dentre as demais, nove delas foram classificadas como médio e outras quatro como fraco, sendo 13, por conseguinte, a quantidade de questões nos dois patamares mais baixos de discriminação.
Constata -se, assim, que a prova – no que se refere ao Componente de Conhecimento Específico – possuía capacidade mediana de discriminar entre aqueles que dominam ou não o conteúdo.
Dentre as questões que alcançaram os maiores índices de discriminação, quatro delas, as de números 12, 18, 20 e 22, foram classificadas com índice de discriminação muito bom, situando -se no intervalo de 0,40 a 0,47 do índice, e todas elas foram classificadas na categoria médio.
A questão de número 32 foi a mais difícil dentre as 26 questões específicas válidas, com baixo índice de facilidade, apenas 8,0 \% de acertos.
Essa questão apresentou poder discriminatório igualmente baixo, 0,06, o que comprova ter sido esta questão a mais difícil para os estudantes.
Destaca -se, também, a questão 34, com índice de facilidade 0,14, o que, em termos percentuais, corresponde a 14,0 \% de estudantes que responderam acertadamente.
Já 0,14 foi o seu índice de discriminação.
Tais questões foram, portanto, pelo critério ponto bisserial, consideradas inadequadas.
Por isso, as questões 32 e 34 foram eliminadas do cômputo da nota final.
Além destas duas, as demais questões com índice de discriminação fraco, questões 16 e 17 também não foram utilizadas no cômputo final das notas, num total de quatro questões eliminadas.
A título de exemplo das análises do comportamento das questões objetivas, o Gráfico 3.5 analisa a questão 18 do Componente de Conhecimento Específico.
Esta questão apresentou índice de facilidade 0,53 ( médio ), ou seja, 53,0 \% dos estudantes assinalaram acertadamente a opção A, correspondente ao gabarito.
Seu índice de discriminação foi igual a 0,47, classificado como muito bom, também sendo esta questão a que apresentou o maior índice de discriminação.
Neste gráfico, cada uma das cinco curvas representa o percentual de respostas em determinada alternativa da questão 18, em função do número de acertos dos estudantes nessa parte da prova, antes de possíveis eliminações de questões pelo critério do ponto bisserial.
A alternativa correta A, representada no gráfico pela curva em azul, foi escolhida em maiores proporções pelos estudantes com desempenho melhor nessa parte da prova.
Já as alternativas incorretas, também denominadas distratores, foram selecionadas, principalmente, por aqueles com notas mais baixas.
Neste caso também a soma não é sempre 100 \% por causa das questões não respondidas ou com mais de uma opção marcada.
Aqueles com nota zero, na sua quase totalidade deixaram esta questão em branco ou marcaram mais de uma alternativa, comportamento considerado inválido.
A proporção de estudantes que selecionou a resposta correta A aumenta gradativamente, chegando a atingir 100 \% para quem obteve de 14 a 17 acertos, enquanto a proporção dos que escolheram alternativas incorretas decai, a partir de dois acertos, como função do número de acertos nessa parte da prova.
A alternativa B ( linha verde ) apresenta um crescimento no final do gráfico, compensado pela queda na linha azul ( alternativa A ).
Esses alunos com 17 e 18 acertos foram poucos, o que não afetou o índice de discriminação bisserial.
Os gráficos relativos às demais questões do Componente de Conhecimento Específico constam do Anexo I. 3.2.3 Componente de Conhecimento Específico – Física ( Licenciatura ) A Tabela 3.11 apresenta as estatísticas básicas em relação às questões objetivas do Componente de Conhecimento Específico da prova de Física ( Licenciatura ), por Grande Região.
A média do Brasil deste componente foi de 34,3.
A menor média foi observada na região Norte ( 28,5 ), e a maior, na região Sudeste ( 39,0 ).
O desvio padrão de todo o Brasil foi 16,2, sendo o menor desvio padrão encontrado na região Norte ( 12,7 ), e o maior, na região Sudeste ( 17,3 ).
A mediana de todo o Brasil foi 31,6, a mesma encontrada na região Nordeste.
As demais regiões apresentaram valores menores, 26,3 na região Norte, e maiores para a mediana, regiões Sudeste e Sul ( 36,8 em ambas ) e 34,2 na região Centro-Oeste.
A nota máxima da prova foi 89,5, obtida nas questões objetivas do Componente de Conhecimento Específico, por, pelo menos, um estudante nas regiões Nordeste, Sudeste e Sul.
Nas demais regiões, a nota máxima da prova foi 84,2 na região Norte, e 78,9 na região Centro-Oeste.
Em todas as regiões a nota mínima foi zero, sem exceção.
A Tabela 3.12 apresenta os índices de facilidade e discriminação ( ponto bisserial ) das questões objetivas do Componente de Conhecimento Específico para os estudantes de Física ( Licenciatura ).
Para facilitar a diferenciação das questões usaram -se as mesmas cores das Tabelas 3.8 e 3.10 para as diferentes classificações dos índices de facilidade e de discriminação.
Dentre as questões objetivas da parte da prova relativa ao Componente de Conhecimento Específico, nenhuma das questões foi anulada pela Comissão.
Desse modo, a classificação, quanto ao índice de facilidade, foi estabelecida com base em todas as 27 questões.
A partir dos índices obtidos, pode -se concluir que mais de três quartos das questões objetivas da prova foi considerado, pelo menos, difícil : das 27 questões, 22 foram classificadas como difícil ( a classificação modal ) ou como muito difícil.
Uma questão foi classificada como fácil, e outras quatro consideradas como médio.
Já quanto aos índices de discriminação das questões objetivas do Componente de Conhecimento Específico da prova, tem -se como resultado a seguinte classificação : sete das 27 questões foram consideradas como boas, enquanto quatro delas tiveram índice de discriminação muito bom.
Assim, para a um pouco mais de 1/3 das questões – 11 em 27 – os índices de discriminação foram bom ou muito bom.
Dentre as demais, oito delas foram classificadas como médio e outras oito como fraco, sendo 16, por conseguinte, a quantidade de questões nos dois patamares mais baixos de discriminação.
Constata -se, assim, que a prova – no que se refere ao Componente de Conhecimento Específico – possuía capacidade fraca de discriminar entre aqueles que dominam ou não o conteúdo.
Dentre as questões que alcançaram os maiores índices de discriminação, quatro delas, as de números 15, 26, 29 e 31, foram classificadas com índice de discriminação muito bom, situando -se no intervalo de 0,40 a 0,45 do índice, uma delas ( questão 29 ) foi classificada na categoria fácil, quanto ao índice de facilidade, duas A Tabela 3.12 apresenta os índices de facilidade e discriminação ( ponto bisserial ) das questões objetivas do Componente de Conhecimento Específico para os estudantes de Física ( Licenciatura ).
Para facilitar a diferenciação das questões usaram -se as mesmas cores das Tabelas 3.8 e 3.10 para as diferentes classificações dos índices de facilidade e de discriminação.
Dentre as questões objetivas da parte da prova relativa ao Componente de Conhecimento Específico, nenhuma das questões foi anulada pela Comissão.
Desse modo, a classificação, quanto ao índice de facilidade, foi estabelecida com base em todas as 27 questões.
A partir dos índices obtidos, pode -se concluir que mais de três quartos das questões objetivas da prova foi considerado, pelo menos, difícil : das 27 questões, 22 foram classificadas como difícil ( a classificação modal ) ou como muito difícil.
Uma questão foi classificada como fácil, e outras quatro consideradas como médio.
Já quanto aos índices de discriminação das questões objetivas do Componente de Conhecimento Específico da prova, tem -se como resultado a seguinte classificação : sete das 27 questões foram consideradas como boas, enquanto quatro delas tiveram índice de discriminação muito bom.
Assim, para a um pouco mais de 1/3 das questões – 11 em 27 – os índices de discriminação foram bom ou muito bom.
Dentre as demais, oito delas foram classificadas como médio e outras oito como fraco, sendo 16, por conseguinte, a quantidade de questões nos dois patamares mais baixos de discriminação.
Constata -se, assim, que a prova – no que se refere ao Componente de Conhecimento Específico – possuía capacidade fraca de discriminar entre aqueles que dominam ou não o conteúdo.
Dentre as questões que alcançaram os maiores índices de discriminação, quatro delas, as de números 15, 26, 29 e 31, foram classificadas com índice de discriminação muito bom, situando -se no intervalo de 0,40 a 0,45 do índice, uma delas ( questão 29 ) foi classificada na categoria fácil, quanto ao índice de facilidade, duas delas ( questões 15, e 31 ), na categoria médio, e a outra ( questão 26 ), na categoria difícil.
A questão de número 11 foi a mais difícil dentre as 27 questões específicas, com baixo índice de facilidade, apenas 10,0 \% de acertos.
Essa questão apresentou poder discriminatório igualmente baixo, 0,09, o que comprova ter sido esta questão a mais difícil para os estudantes.
Destaca -se, também, a questão 19, com índice de facilidade 0,12, o que, em termos percentuais, corresponde a 12,0 \% de estudantes que responderam acertadamente.
Já 0,15 foi o seu índice de discriminação.
Tais questões foram, portanto, pelo critério ponto bisserial, consideradas inadequadas.
Por isso, as questões 11 e 19 foram eliminadas do cômputo da nota final.
Além destas duas, as demais questões com índice de discriminação fraco, questões 13, 16, 17, 21, 30 e 33 também não foram utilizadas no cômputo final das notas, num total de oito questões eliminadas.
A título de exemplo das análises do comportamento das questões objetivas, o Gráfico 3.6 analisa a questão 26 do Componente de Conhecimento Específico.
Esta questão foi classificada como difícil.
Apresentou índice de facilidade 0,32, ou seja, 32,0 \% dos estudantes assinalaram acertadamente a opção A, correspondente ao gabarito.
Seu índice de discriminação foi igual a 0,45, classificado como muito bom, também sendo esta questão a que apresentou o maior índice discriminatório.
Neste gráfico, cada uma das cinco curvas representa o percentual de respostas em determinada alternativa da questão 26, em função do número de acertos dos estudantes nessa parte da prova, antes de possíveis eliminações de questões pelo critério do ponto bisserial.
A alternativa correta A, representada no gráfico pela curva em azul, foi escolhida em maiores proporções pelos estudantes com desempenho melhor nessa parte da prova.
Já as alternativas incorretas, também denominadas distratores, foram selecionadas, principalmente, por aqueles com notas mais baixas.
Neste caso também a soma não é sempre 100 \% por causa das questões não respondidas ou com mais de uma opção marcada.
Aqueles com nota zero, na sua quase totalidade deixaram esta questão em branco ou marcaram mais de uma alternativa, comportamento considerado inválido.
A proporção de estudantes que selecionou a resposta correta A aumenta gradativamente ( com algumas oscilações ), chegando a atingir 100 \% para 20 acertos, enquanto a proporção dos que escolheram alternativas incorretas decai, a partir do primeiro acerto, como função do número de acertos nessa parte da prova.
Os gráficos relativos às demais questões do Componente de Conhecimento Específico constam do Anexo I. 3.3 ANÁLISE DAS QUESTÕES DISCURSIVAS Esta seção apresenta estatísticas selecionadas e histogramas das Questões Discursivas de Formação Geral ( 3.3.1 ) e Conhecimento Específico ( 3.3.2 ).
São também apresentadas e comparadas as médias de subpopulações caracterizadas por Grande Região. 3.3.1 Componente de Formação Geral As análises dos resultados de desempenho dos estudantes de Física ( Bacharelado e Licenciatura ) nas duas questões discursivas relativas à Formação Geral, considerando -se as notas de conteúdo e de desempenho linguístico, encontram -se na Tabela 3.13 e no Gráfico 3.7.
Na Tabela 3.13, observa -se que a nota média nesse conjunto de questões foi próxima da obtida nas objetivas.
Os estudantes, de todo o Brasil, obtiveram, em Formação Geral, média 58,9 nas questões objetivas e 50,9 nas questões discursivas.
No entanto, pode -se notar um aumento do desvio padrão de 22,3 nas questões objetivas do Componente de Formação Geral dos estudantes de todo o Brasil, para 26,8 nas questões discursivas do mesmo componente.
A maior média foi obtida na região Sul ( 52,7 ), e a menor, na região Norte ( 48,8 ).
A mediana de todo o Brasil, neste componente, foi 57,0, nas regiões Sudeste, Sul e Centro-Oeste a mediana foi maior, 58,5, 59,5 e 58,8, respectivamente.
Nas demais regiões foi menor : Norte ( 53,5 ) e Nordeste ( 56,0 ).
A nota máxima ( 97,5 ) foi obtida por, pelo menos, um estudante na região Nordeste.
Nas demais regiões a nota máxima foi : 94,0 na região Norte, 96,5 na região Sudeste, 95,5 na região Sul e 96,0 na região Centro-Oeste.
A nota mínima ( 0,0 ) foi a mesma em todas as regiões do Brasil, sem exceção.
O Gráfico 3.7 representa a distribuição das notas nas questões discursivas do Componente de Formação Geral, segundo a opção de Bacharelado/Licenciatura.
As modas destas distribuições ocorrem no intervalo [ 0 ; 10 ] para o Bacharelado e no intervalo ( 60:70 ] para Licenciatura.
Para as duas habilitações observam -se outros máximos locais nos intervalos [ 0 ; 10 ] e ( 60 ; 70 ], para Licenciatura e Bacharelado, respectivamente, sendo que no primeiro intervalo incluem -se além das notas zero, a frequência de estudantes que deixaram este tipo de questão em branco.
Na sequência, os resultados verificados para cada uma das questões discursivas de Formação Geral serão apresentados, estabelecendo -se relações com os conteúdos abordados em cada uma delas.
Os comentários da Banca de docentes corretores, a respeito do observado na correção das respostas dos estudantes, suas impressões e conclusões serão apresentados junto à análise de cada questão.
Cumpre esclarecer que, tendo em vista que as questões discursivas de Formação Geral são padronizadas, ou seja, constam de todas as provas, os comentários da Banca são os mesmos para todas as carreiras acadêmicas, sendo direcionados a todos os estudantes que participaram do ENADE/2014.
A seguir, serão analisados os desempenhos da Área de Física ( Bacharelado e Licenciatura ) nas duas questões discursivas de Formação Geral do ENADE/2014, comparando -se os resultados obtidos com comentários para cada questão. 3.3.1.1 Análise de conteúdo da Questão Discursiva 1 do Componente de Formação Geral Os dados de Física ( Bacharelado e Licenciatura ), obtidos a partir das respostas à questão 1, encontram -se na Tabela 3.14 e no Gráfico 3.8.
Nessa questão – de melhor desempenho dentre as duas de Formação Geral – os estudantes, de todo o Brasil, tiveram média 52,5.
A maior média para a questão 1 foi obtida na região Sul ( 55,8 ), e a menor, na região Norte ( 49,0 ).
Quanto à variabilidade das notas, o desvio padrão de todo o Brasil foi 32,6.
O menor desvio padrão foi obtido na região Norte ( 30,9 ), e o maior desvio padrão foi obtido na região Sudeste ( 33,0 ).
As medianas das regiões Nordeste, Centro-Oeste e do Brasil, como um todo, foram iguais ( 55,0 ).
Nas regiões Sudeste e Sul, a mediana foi 60,0 ; e na região Norte, foi 50,0.
As notas máximas e mínimas da questão discursiva 1 foram as mesmas para todas as regiões do Brasil, respectivamente, 100,0 e 0,0.
O Gráfico 3.8 mostra a distribuição das notas na questão discursiva 1 do Componente de Formação Geral para estudantes do Bacharelado e da Licenciatura.
Observa -se que para os de Bacharelado, a maior frequência corresponde aos estudantes que deixaram a questão em branco, e para os de Licenciatura aos que obtiveram notas no intervalo ( 70 ; 80 ]. 3.3.1.2 Comentários sobre a correção de Conteúdo das respostas à Questão Discursiva 1 O enunciado apresentado da Questão 1 era claro e foi redigido em linguagem acessível.
A questão abordava o tema da mobilidade urbana, relacionando -a aos ideais de desenvolvimento urbano sustentável estabelecidos pela ONU, do ponto de vista ecológico, cultural, político, institucional, social e econômico.
O enunciado apresentava uma tabela, sem referência à fonte de publicação, que detalhava percentualmente as modalidades de deslocamento da população urbana brasileira, com base em duas distinções : motorizado versus não motorizado ; coletivo versus individual.
O comando da questão solicitava, então, que o estudante desenvolvesse um texto dissertativo composto de duas partes : análise das consequências do transporte motorizado para o desenvolvimento sustentável e elaboração de duas propostas de intervenção pautadas pelo incentivo ao uso da bicicleta, meio de transporte não motorizado e, portanto, não poluente.
A temática dos efeitos da industrialização, com foco na produção de automóveis, e consequências para o meio ambiente é atual.
Além disso, a questão provoca a reflexão sobre a necessidade de adequação das cidades ao número de habitantes, no contexto da utilização dos meios de transportes de massa : metroviário, ferroviário, por ônibus.
Esperava -se, assim, que o concluinte recorresse a diferentes argumentos para analisar as consequências do transporte motorizado, explicitados pelo texto motivador, e propusesse duas ações de intervenção, as quais considerasse relevantes para o incremento ao uso de bicicleta, que, preferencialmente, deveriam manter coerência com a argumentação desenvolvida na primeira parte da resposta.
O conflito entre a busca pelo transporte individual e o coletivo é inerente a todas as grandes cidades do mundo.
No Brasil, esse conflito emerge com maior gravidade, como efeito do aumento da quantidade de veículos automotores ( automóveis e motocicletas ) e da baixa qualidade dos transportes coletivos.
Tal associação incentiva o uso de meios de transporte individual em detrimento da utilização dos coletivos.
O padrão de respostas era tecnicamente adequado ao enunciado e amplo o suficiente para abranger as variações de respostas apresentadas.
Deve ser ressaltado que a proteção ao ambiente ecologicamente equilibrado é tema conhecido dos concluintes dos cursos universitários, que deveriam possuir informações suficientes para desenvolver uma resposta abordando pelo menos dois dos aspectos listados no padrão de respostas como possibilidades de consequências ( item ‘ a ’ ) e duas propostas de ação de intervenção aceitáveis ( item ‘ b ’ ).
Segundo o padrão, concluintes poderiam desenvolver, quanto ao item ‘ a ’, os seguintes pontos : aumento da emissão de poluentes atmosféricos ; aumento da emissão de gases de efeito estufa ( CO2- dióxido de carbono, CO - monóxido de carbono, O3 - Ozônio ) ; aumento da poluição visual e sonora ; aumento da temperatura local e global ; aumento do consumo de combustíveis ; aumento de problemas de saúde ( cardíacos, respiratórios, dermatológicos ) ; aumento da frota de veículos, promovendo congestionamentos urbanos ; diminuição de áreas verdes ; desmatamento ; aumento das áreas de impermeabilização, resultando em enchentes, diminuição da infiltração da água e recarga de lençóis freáticos ; elevação dos custos de manutenção das cidades ( metroferrovias, rodovias, tratamento de água, limpeza da cidade, etc. ) ; necessidade de ampliação de vias trafegáveis ; necessidade de ampliação de áreas de estacionamento.
Em relação ao item ‘ b ’, os estudantes poderiam considerar os seguintes aspectos : construção de vias exclusivas para bicicletas ( ciclovias e ciclofaixas ) ; proposição de formas de integração entre o transporte por bicicletas, o metroviário e os ônibus coletivos, a fim de garantir segurança e conforto em momentos de adversidades climáticas e relevo acidentado ; pontos de aluguel e/ou empréstimos de bicicletas ; construção de bicicletários ; investimentos na segurança pública ; políticas de incentivo ao uso de bicicletas ( educação ambiental, qualidade de vida, saúde, propaganda ) ; implementação de políticas de crédito e de redução de custos das bicicletas.
Para o item ‘ a ’, foi encontrado um número expressivo de respostas adequadas.
As consequências mais citadas estavam relacionadas com a emissão de poluentes e o aumento dos engarrafamentos.
Houve críticas frequentes à qualidade do transporte público e à necessidade de locomoção rápida diante da distância das residências aos locais de trabalho.
Em tom de reivindicação, muitos estudantes defenderam solução para os problemas do transporte coletivo, quer metroviário, ferroviário ou por ônibus.
O uso da bicicleta também foi considerado por muitos como um meio de transporte benéfico à saúde, mas alguns apontaram a impossibilidade de sua utilização pelos idosos e deficientes.
No concernente ao item ‘ b ’, de modo quase uniforme citou -se a necessidade de ciclovias e ciclofaixas, bem como de se garantir segurança para os ciclistas.
A necessidade de manutenção das ciclovias e ciclofaixas existentes em algumas cidades foram também lembradas.
Com menos frequência também houve a indicação de maior conscientização da população, inclusive pela mídia, e pelo incentivo à aquisição desses veículos não motorizados, por meio da diminuição de impostos para redução do custo de aquisição.
Houve quem sugerisse, inclusive, uma política de doação de bicicletas.
Alguns mencionaram a necessidade de se assegurar integração das ciclovias com os demais modais, com a construção de bicicletários seguros diante das grandes distâncias a percorrer, notadamente nas metrópoles.
Também com base na dificuldade de o veículo não motorizado vencer as longas distâncias que o trabalhador deve percorrer nos centros urbanos para chegar ao local onde exerce a sua atividade produtiva, alguns estudantes argumentaram contra a utilização da bicicleta como forma de resolução da poluição ambiental. bicicletas, como se requeria no comando da questão.
Nesse sentido, foram encontradas respostas que indicavam a caminhada como opção.
Foram consideradas respostas fracas ( notas de zero a 30 ) aquelas meramente opinativas, ou que apenas transcreveram elementos do enunciado, sem qualquer acréscimo ou apreciação crítica.
As respostas medianas ( notas entre 35 e 70 ) foram as que não apresentavam algum desenvolvimento para os tópicos citados ( duas consequências e duas ações de intervenção ) ou por falta de alguns deles.
Já as respostas boas ( notas de 75 a 100 ) continham os tópicos considerados pelo padrão de respostas, com algum desenvolvimento, pelo menos, regular de argumentação.
As respostas dos estudantes revelaram um adequado conhecimento quanto aos efeitos do transporte motorizado para o ambiente e que a temática da ecologia parece ser tema bem compreendido, aparecendo nas respostas como uma preocupação para as futuras gerações.
Nesse sentido, algumas políticas públicas foram mencionadas apropriadamente para assegurar o desenvolvimento sustentável pelos respondentes, tais como : educação, habitação, saneamento, saúde, transporte e segurança.
No entanto, de modo geral, as respostas demostraram que os estudantes têm muita dificuldade na expressão escrita do pensamento, como se pode constatar pela avaliação do desempenho linguístico que ficou a cargo de uma banca específica, formada por profissionais da área de Língua Portuguesa. 3.3.1.3 Análise de Conteúdo da Questão Discursiva 2 do Componente de Formação Geral A Tabela 3.15 mostra que o desempenho médio dos estudantes na questão discursiva 2 ( média 48,6 ) foi inferior ao obtido na questão discursiva 1 ( média 52,5 ).
A região Centro-Oeste foi aquela cuja média, nessa questão, foi maior ( 49,9 ), e a de menor média foi a região Norte ( 46,8 ).
Quanto à variabilidade das notas, o desvio padrão de todo o Brasil foi 33,4, superior ao obtido na questão discursiva 1 ( 32,6 ).
O maior desvio nessa questão foi obtido na região Centro-Oeste ( 34,8 ), enquanto o menor foi obtido na região Norte ( 32,3 ).
A mediana de todo o Brasil foi 60,0, a mesma das regiões Nordeste, Sudeste e Sul, enquanto a menor mediana foi encontrada nas regiões Norte e Centro-Oeste ( 50,0 ).
As notas máximas ( 100,0 ) foram as mesmas em todas as regiões do Brasil, assim como as notas mínimas ( 0,0 ), sem exceção.
O Gráfico 3.9 mostra a distribuição das notas na questão discursiva 2 do Componente de Formação Geral para estudantes do Bacharelado e da Licenciatura.
Observa -se que a maior frequência corresponde aos estudantes que deixaram a questão em branco tanto para aqueles de Bacharelado quanto para os de Licenciatura. 3.3.1.4 Comentários sobre a correção de Conteúdo das respostas à Questão Discursiva 2 A questão 2 abordava o tema da violência urbana a partir de um texto motivador recolhido em um site de um jornal brasileiro.
O fragmento relatava a prisão de três jovens por agressão a um outro jovem, menor de idade, guardador de carros.
A notícia trazia informações sobre o estado do jovem agredido e sobre as possíveis penas a serem aplicadas aos agressores.
O enunciado solicitava, então, que o estudante desenvolvesse um texto dissertativo composto de duas partes : análise de duas causas do tipo de violência descrito no texto e apresentação de dois fatores que contribuiriam para evitar a agressão relatada.
Como o texto motivador reduzia -se ao relato de um episódio de violência urbana, o enunciado não oferecia argumentos que ajudassem o estudante a elaborar uma análise das possíveis causas do fato.
As propostas a serem apresentadas, no sentido de se evitar em situações como a relatada, deveriam manter coerência com a análise das causas, desenvolvida na primeira parte da resposta.
O enunciado era claro e a matéria permitia identificar tratar -se de um conflito tipicamente urbano no qual se envolveram jovens de baixa renda em disputa por um território onde exerciam trabalho informal.
Sendo essa uma das hipóteses albergada sobre o amplo tema da violência urbana no qual se desenvolve o relato.
A temática da questão é bastante conhecida, não somente por debates no âmbito universitário, especialmente na área de Ciências Sociais, bem como pela ênfase que é dada pela mídia onde assume contornos, muitas vezes, maiores do que o problema em si.
Assim, por tratar -se de tema bastante presente e debatido no cotidiano das pessoas, a questão pode ser considerada de baixa complexidade.
O padrão de resposta apresentava uma gama de possibilidades de causas da violência aceitas como corretas para o item ‘ a ’, tais como : problemas relacionados à educação ( baixa escolaridade, evasão escolar, qualidade da educação, distanciamento entre escola e a realidade social, tempo de permanência na escola ) ; desigualdades socioculturais ( gênero, etnia, economia, etc. ) ; desemprego e falta de qualificação profissional ; precariedade da segurança pública ; uso de drogas ; desvalorização da vida humana ; banalização da violência ; sensação de impunidade ; ausência de políticas sociais ; degradação da vida urbana ; desconhecimento ou desrespeito aos direitos humanos e constitucionais ; desestruturação familiar ; desvalorização de princípios éticos e morais.
No concernente ao item ‘ b ’, o padrão também previa a possibilidade de apresentação de diversos fatores que podem contribuir para evitar a violência : políticas de segurança mais efetivas ; políticas públicas de melhoria das condições socioeconômicas ; maior consciência cidadã e respeito à vida ; melhor distribuição de renda ; melhoria da educação ( aumento da escolaridade, redução da evasão escolar, qualidade da educação, aproximação entre a escola e a realidade social, aumento do tempo de permanência na escola ) ; aumento da oferta de emprego e melhoria da qualificação profissional ; medidas preventivas ao uso de drogas ; maior eficácia do sistema judiciário ; revisão da legislação penal ; valorização de princípios éticos, morais e familiares.
O enunciado da questão permitia que os estudantes, em suas respostas, pautassem o problema por diversos ângulos.
Dessa forma, o padrão buscou indicar diversas possibilidades a serem identificadas pelo corretor na leitura da resposta do estudante.
As notas foram atribuídas considerando o desenvolvimento da resposta e a indicação de ao menos dois dos itens exigidos pelo comando da questão nos itens ‘ a ’ ( duas causas ) e ‘ b ’ ( dois fatores ), dentre os considerados pelo padrão de respostas.
No entanto, o padrão exigia que as respostas estivessem afinadas com o perfil profissional previsto no Art. 3º da Portaria nº 255, de 02 de junho de 2014, onde se lê : Art. 3º No componente de Formação Geral serão considerados os seguintes elementos integrantes do perfil profissional : atitude ética ; comprometimento social ; compreensão de temas que transcendam ao ambiente próprio de sua formação, relevantes para a realidade social ; espírito científico, humanístico e reflexivo ; capacidade de análise crítica e integradora da realidade ; e aptidão para socializar conhecimentos em vários contextos e públicos diferenciados.
Nesse sentido, teses repressivas não conformadas com o texto constitucional e amparados pelos direitos humanos, não foram consideradas.
Questões como a majoração das penas, instituição de penas mais duras e criminalização da conduta dos menores desviados, só foram consideradas no caso de serem tratadas como menção a um debate que está presente, nunca como medida a ser aplicada de forma arbitrária e inconstitucional.
Também não foram previstas visões mais autoritárias e repressoras como a defesa de proibição do trabalho informal, como é o caso dos “ flanelinhas ”, e a defesa de não fornecer dinheiro aos “ guardadores ” de carro, identificado como esmolas por alguns.
Não foram consideradas respostas que meramente reproduziam o enunciado sem qualquer reflexão ou análise.
Aquelas que utilizaram os elementos do enunciado, mesmo que apenas transcritos, mas desenvolveram algum raciocínio pessoal ou indicaram causas e/ou fatores foram valoradas.
Dentre as respostas apresentadas, a quantidade de provas boas foi expressiva.
Nessa questão, 14 \% dos estudantes de exame deixaram o espaço para resposta em branco.
Entre as respostas corrigidas em torno de 3 \% receberam nota zero, e a média das notas foi 5,9 pontos.
As causas apontadas com maior frequência para a violência foram : a ausência de educação e o desemprego, com soluções correlatas de educação profissional e de qualidade e abertura de postos de trabalho.
A questão das drogas também foi bastante mencionada, e a solução correlata mais indicada foi o acolhimento de usuários de drogas.
Outro tema bastante citado foi a dissolução das famílias e dos valores familiares.
Tendo em vista que o texto falava de moradores de rua, a ausência de moradia apareceu em muitas respostas, bem como a disputa territorial e a necessidade de sobrevivência.
Nesse sentido, as propostas apontavam para a necessidade de ajuda da área de Serviço Social das prefeituras e atendimento psicológico gratuito.
Numa ótica mais repressora, surgiram críticas quanto à leniência das leis, acompanhadas da indicação da necessidade de atuação mais rigorosa do Poder Judiciário, como também da aplicação de penas maiores, negativa de esmolas, proibição do trabalho informal ( guardador de carros ), dentre outras.
Algumas respostas indicaram como causa o modelo capitalista, o consumismo e a ganância.
Outras indicaram o trabalho infantil e a correlata necessidade de aplicação do ECA ( Estatuto da Criança e do Adolescente ).
Em relação ao tema violência, inúmeras respostas apresentaram visões de intolerância quanto aos menores infratores, postulando a diminuição da idade para efeitos de criminalização, bem como defendendo a majoração de penas.
Essa perspectiva revela uma cisão social, de certa forma estimulada pelo tipo de divulgação da mídia quanto aos temas em foco, o que, sem dúvida, ficou evidenciado em algumas respostas.
A existência de inúmeras respostas pugnando pela maior repressão estatal e pelo aumento das penas, acrescido da busca pela diminuição da idade penal, revela que muitos estudantes não foram atingidos por uma formação superior que se exige humanista, ética e comprometida socialmente, que busque os fundamentos para resolver o problema da violência por meio de práticas democráticas e de inclusão social.
As respostas desse tipo também revelaram um alto grau de intolerância para com os moradores de rua, associando -os muitas vezes ao tráfico de drogas, à exploração indevida do espaço público, bem como a crimes como extorsão dos motoristas quando pedem dinheiro para guardar os veículos.
Nessa trilha, verifica -se lacuna na discussão de temas importantes para a formação de profissionais de nível superior tais como : sociodiversidade, multiculturalismo e violência ; tolerância/intolerância ; inclusão/exclusão.
Para aqueles que responderam em conformidade com o padrão de resposta verificou -se uma plena formação social, conclamando pelo auxílio não só do Estado, mas também da Sociedade no amparo daqueles que não têm as mínimas condições de sobrevivência e que estão nas ruas. 3.3.1.5 Análise de Língua Portuguesa das Questões Discursivas do Componente de Formação Geral Os dados de Física, obtidos a partir das respostas às questões discursivas do Componente de Formação Geral no que tange à Língua Portuguesa, encontram -se na Tabela 3.16 e no Gráfico 3.10.
Nesse aspecto, os estudantes, de todo o Brasil, tiveram média 52,5.
A maior média com respeito à Língua Portuguesa foi obtida na região Sul ( 54,2 ), e a menor, nas regiões Sudeste e Centro-Oeste ( 52,2 em ambas ).
Quanto à variabilidade das notas, o desvio padrão de todo o Brasil foi 24,9.
O menor desvio padrão foi obtido na região Norte ( 23,9 ) e o maior desvio padrão foi obtido na região Sudeste ( 26,0 ).
As medianas das notas de Língua Portuguesa para as regiões Norte, Nordeste e o Brasil, como um todo, foram 60,0, enquanto nas demais regiões a mediana foi 62,5.
A nota máxima para todo o Brasil foi de 92,5, com, pelo menos, um estudante tirando essa nota na região Nordeste, enquanto nas demais regiões a nota máxima foi 90,0.
Além disso, a nota mínima foi zero em todas as regiões do país, sem exceção.
O Gráfico 3.10 mostra a distribuição das notas de Língua Portuguesa do Componente de Formação Geral.
Observa -se que, para ambas as habilitações, a maior frequência corresponde aos estudantes que obtiveram nota no intervalo ( 60 ; 70 ]. 3.3.1.6 Comentários sobre a correção das respostas de Formação Geral com respeito à Língua Portuguesa Ao encaminhar as questões 1 e 2 na direção da produção de um texto dissertativo, esperava -se que o estudante utilizasse seus conhecimentos sobre o assunto e estruturasse seus textos de acordo com as exigências do registro formal próprio dessa situação comunicativa.
Essa configuração determina exigências quanto : à adequação da seleção vocabular, ao desenvolvimento do conteúdo, à estruturação sintática dos períodos, à organização lógica das ideias, à utilização de procedimentos de encadeamento textual e referenciação, à obediência às exigências morfossintáticas próprias da modalidade escrita da norma-padrão, ao respeito às regras ortográficas e às regras de acentuação gráfica.
O padrão de respostas utilizado na avaliação considerou os aspectos relevantes ao bom desempenho linguístico como competências distintas, de modo a permitir um mapeamento detalhado do domínio dos recursos disponíveis na Língua Portuguesa para a comunicação escrita formal.
Com base nesse objetivo, foram avaliados os seguintes aspectos : a ) Estruturação textual condizente com o gênero solicitado e o modo de organização textual expositivo adequado ao gênero – essa competência envolve : a estruturação sintática condizente com o padrão da modalidade escrita formal da língua portuguesa, de modo a garantir a clareza necessária ; a distribuição do conteúdo do texto em parágrafos, de modo a garantir a sua organização temática ; a utilização de operadores discursivos que contribuam para a progressão temática do texto, estabelecendo relações lógicas entre as ideias apresentadas, tanto do ponto de vista intrafrasal, como do interfrasal ; a utilização de procedimentos de referenciação lexical e pronominal que permitam a retomada de referentes textuais ; o respeito às regras de pontuação como fator de estruturação do período.
Espera -se, portanto, que o estudante recorra a procedimentos linguístico-discursivos para organizar seu texto, permitindo o encadeamento lógico entre suas partes de forma a garantir a progressão e a coerência textuais.
Isso significa que os seguintes procedimentos foram penalizados, de acordo com o padrão de respostas proposto : • a estrutura lógico-gramatical do texto fica comprometida pela elaboração de frases fragmentadas ; • sequência justaposta de ideias sem encaixamentos sintáticos, reproduzindo hábitos da oralidade ; • elaboração de frase com apenas oração subordinada, sem oração principal ; • emprego equivocado do conector ( preposição, conjunção, pronome relativo, alguns advérbios e locuções adverbiais ) que não estabeleça relação lógica entre dois trechos do texto e prejudique a compreensão da mensagem ; • emprego do pronome relativo sem a preposição, quando obrigatória ; • repetição ou substituição inadequada de palavras sem utilização dos recursos oferecidos pela língua ( pronome, advérbio, artigo, sinônimo ) ;  emprego inadequado dos pronomes relativos “ cujo(a) ” e “ onde ” ;  utilização inadequada dos sinais de pontuação que comprometa a clareza textual. b ) Respeito às convenções ortográficas da norma-padrão da Língua Portuguesa – essa competência envolve o domínio das regras de acentuação gráfica e da grafia padrão das palavras ( com ausência de abreviaturas próprias da linguagem da internet ), de acordo com as convenções estabelecidas pela legislação em vigor e consubstanciadas no Vocabulário Ortográfico da Língua Portuguesa, editado pela Academia Brasileira de Letras ( com aceitação da legislação anterior, no caso das regras relativas ao uso do hífen e da acentuação gráfica ).
Espera -se que o participante :  grafe corretamente as palavras ;  respeite as regras de acentuação gráfica ;  empregue maiúsculas em início de frase, em nomes próprios de pessoas, lugares ou instituições ;  evite abreviações como p/, vc, tb, pq, tá, né, usadas muitas vezes em escrita informal e na internet ;  obedeça às regras de separação de sílabas no final da linha. c ) Domínio dos diferentes aspectos morfossintáticos próprios da modalidade escrita formal da norma-padrão da Língua Portuguesa – essa competência envolve : a concordância nominal, a concordância verbal, a regência nominal, a regência verbal, a flexão nominal, a flexão verbal, a correlação entre os tempos verbais, a colocação pronominal e a utilização de sinais de pontuação que contribuam para a organização lógica da frase e do texto.
Espera -se que o participante :  flexione o verbo para estabelecer concordância de número com o sujeito da frase ;  flexione o artigo, o adjetivo e o pronome para concordar em número e em gênero com o substantivo a que se referem ;  observe a regência nominal e a verbal, utilizando a preposição adequada depois de um substantivo, um verbo ou um adjetivo ;  empregue adequadamente o acento grave indicador de crase nos casos em que se fizer necessário ;  obedeça às regras de colocação pronominal ( próclise e ênclise ), distintas dos hábitos da oralidade ou da escrita informal ; d ) Seleção vocabular adequada à modalidade escrita formal da Língua Portuguesa, exigida pela situação comunicativa – essa competência envolve a precisão na utilização do vocabulário relacionado à temática solicitada pela questão ; a ausência de marcas da oralidade, como termos de sentido de muito genérico ( “ coisa ”, “ negócio ”, “ você ” ) e termos de registros mais informais ( como gírias, jargões, frases feitas, ditados populares, termos regionais ).
Assim, espera -se que o participante respeite a adequação vocabular não utilizando gírias ou expressões coloquiais, evite repetição desnecessária de palavras e utilize um vocabulário mais formal, adequado ao texto de caráter dissertativo.
A escolha dessas competências para subsidiar o processo de avaliação apoia -se na concepção de que, no desempenho dos graduandos, a modalidade escrita tem apresentado uma intensa simplificação, originada no padrão da modalidade oral da Língua Portuguesa.
No caso do texto de base dissertativa, inscrito em um registro formal, a distância entre as duas modalidades é ainda maior, o que provoca situações de hipercorreção ( desvios provocados pela incorporação indevida de uma regra da norma-padrão ) e de truncamentos sintáticos ( estruturas frasais incompreensíveis devido à complexidade sintática própria da modalidade escrita ).
Observam -se, então, os seguintes aspectos que marcam essa distinção entre as duas modalidades, devido à excessiva simplificação da modalidade falada : a ) redução drástica de estruturas subordinadas, compensada pelo aumento na frequência de estruturas coordenadas e absolutas, por um lado, ou pela elaboração de estruturas truncadas pelo excesso de ideias sem a devida conexão subordinativa ; b ) redução no uso de conectores para expressar relações lógicas essenciais à construção do texto, substituídas pela exigência de inferência por parte do interlocutor para suprir a sua ausência ; c ) redução cada vez maior do uso do subjuntivo, ao lado da ampliação do uso do indicativo combinado a estruturas frasais coordenadas ou absolutas ; d ) empobrecimento do processo de referenciação, com a repetição exaustiva de pronomes ou nomes ; e ) simplificação extrema da marcação da categoria tempo na morfologia verbal ; f ) falta de domínio de vocabulário mais abstrato e de maior complexidade, essencial ao desenvolvimento do processo dissertativo ; g ) redução drástica no emprego da acentuação gráfica, processo intensificado pela divulgação imprecisa das mudanças promovidas pelo último acordo ortográfico.
Os aspectos macroestruturais da elaboração do texto não foram avaliados neste processo, para não penalizar duplamente os estudantes, já que a banca de formação geral, composta por profissionais de diferentes áreas do conhecimento, encarregou -se da avaliação do conteúdo desenvolvido nas questões.
São eles : progressão temática, coerência na relação com os conhecimentos de mundo dos usuários da língua, inteligibilidade, atendimento ao solicitado no enunciado do ponto de vista do desenvolvimento do conteúdo, entre outros.
A grade de avaliação do desempenho linguístico considerou, portanto, três grandes grupos de competências, segundo os aspectos explicitados anteriormente : 1.
Aspectos ortográficos : domínio das convenções ortográficas : grafia de vogais e consoantes, uso de maiúsculas e minúsculas, emprego do hífen e acentuação gráfica ; 2.
Aspectos textuais : domínio dos procedimentos de estruturação textual do ponto de vista microestrutural : organização interna dos períodos, emprego de conectores para a articulação lógica entre os períodos e entre os parágrafos, emprego de marcas de referenciação lexical e pronominal ; utilização dos sinais de pontuação que contribuem para a organização lógica da frase. 3.
Aspectos morfossintáticos e vocabulares : domínio das regras de caráter morfossintático estabelecidas como modelares do ponto de vista da modalidade escrita formal da norma-padrão da Língua Portuguesa : concordância nominal e verbal, regência nominal e verbal, colocação pronominal, flexão nominal e verbal, correlação entre tempos e modos verbais, ausência de marcas de oralidade.
A seleção vocabular adequada à modalidade escrita formal da Língua Portuguesa foi incorporada a essa última competência, tendo em vista a intersecção entre as duas do ponto de vista das exigências do registro formal da modalidade escrita da norma-padrão.
Os resultados da avaliação correspondem aos seguintes aspectos observados em cada competência : Aspectos ortográficos : A correção foi realizada classificando os textos em cinco níveis, nível zero a quatro.
O desempenho dos estudantes ficou concentrado, majoritariamente, no nível 3 ( de 4 a 7 desvios ), devido, principalmente, aos desvios de acentuação.
O número de desvios de grafia foi reduzido.
Diferentemente do resultado do ENADE/2013, aumentou a porcentagem de textos sem desvios ( enquadrados no nível 4 ) e diminuiu a porcentagem de textos enquadrados no nível 2.
Alguns textos, com pior desempenho, foram enquadrados no nível 1 ( de 8 a 12 desvios ), enquanto o nível zero foi atribuído a pouquíssimos casos, já que ele revela falta absoluta de domínio das convenções ortográficas.
Observou -se, portanto, que existe uma diferença muito grande de desempenho nos dois aspectos analisados : baixo índice de desvios ortográficos e grande índice de desvios de acentuação.
Em vários casos, ocorre ausência completa de acentuação gráfica.
Os resultados revelam, portanto, que a tendência dominante entre os universitários brasileiros é a eliminação da acentuação gráfica, provavelmente motivada pela vivência dos jovens relacionada aos aplicativos de comunicação via internet ( redes sociais e e-mails ).
Nesse tipo de comunicação, devido ao ritmo intenso de troca de mensagens, o uso de acentos gráficos foi praticamente abolido.
Outro fator que pode ter relação com essa tendência é a ausência de esclarecimento dos meios de comunicação, das autoridades e das escolas sobre as decisões do Acordo Ortográfico da Língua Portuguesa de 1990, gerando um estado de indefinição para os estudantes.
Os casos mais sistemáticos de eliminação do acento indicador da sílaba tônica são : • palavras proparoxítonas ( “ ridiculos ”, “ publicas ”, “ lideres ”, “ politicos ” ) ; • palavras paroxítonas terminadas em ditongo crescente ( “ homicidios ”, “ latrocinio ”, “ individuo ”, “ dependencia ” ) ; • palavras oxítonas ( “ ninguem ”, “ esta ”, “ ate ”, “ ai ” ).
Por outro lado, destaca -se o uso indevido do acento gráfico em determinadas palavras, como observado nas grafias * “ jornáis ”, * “ telejornaís ”, * “ propíciar ”, * “ medídas ”, * “ dígnidade ”, * “ cídades ”.
Quanto ao domínio das convenções relativas à grafia das palavras, observam -se desvios como : a hipercorreção pela escolha de “ e ” no lugar de “ i ”, por influência de hábitos da oralidade ( * “ enumeros ” por “ inúmeros ”, * “ entevenção ” por “ intervenção ” ) ; a eliminação do “ r ” marcador do infinitivo verbal ( * “ esta ” no lugar de “ estar ” ).
Outros casos de desvios de grafia relacionados à variação diastrática podem ser observados em *“estrupo ”, *“automovís ”, *“viensse ”, *“camihada ”.
São muito frequentes os seguintes desvios de caráter ortográfico, com repercussão morfossintática : • eliminação da marca de infinitivo ( -r- ) e substituição por acento agudo ( “ está ” no lugar de “ estar ” ) ; • confusão entre “ ão ” e “ am ” nas formas verbais ( “ invadão ” no lugar de “ invadam ” e “ estam ” no lugar de “ estão ” no presente do indicativo ; “ estaram ” no lugar de “ estarão ” no futuro do indicativo ) ;  confusão entre a grafia do verbo “ haver ” ( “ há ” ) e o artigo definido ou a preposição “ a ” ;  uso de hífen para separar pronome átono – tanto uso indevido quanto omissão ( no pretérito imperfeito e futuro do subjuntivo : “ evitar -mos ” no lugar de “ evitarmos ”, “ percebesse ” no lugar de “ percebe -se ” e vice-versa ) ; Observam -se, também, muitos casos de inadequação no uso da maiúscula : ausência de diferença entre a primeira letra e as outras, em início de período, principalmente dos estudantes que adotam a escrita em letra de imprensa ; utilização de maiúscula para destacar determinadas palavras-chave do texto, como “ Violência ”, “ Brasileiros ”, “ Fatores ”, “ Ozônio ”, “ Sustentável ”.
Vale observar, também, que, ao contrário do que se esperava, não apareceram abreviaturas próprias do “ internetês ”, ou seja, dos hábitos de comunicação escrita adquiridos pelo uso de redes sociais e emails.
Aspectos textuais : Quanto a esses aspectos a correção também classificou os textos em cinco níveis ( zero a quatro ), em função da quantidade de erros apresentados.
O desempenho dos estudantes ficou concentrado, majoritariamente, nos níveis 3 e 2, devido à grande ocorrência de problemas de estruturação textual.
Foi muito baixo o número de textos que não apresentaram qualquer problema estrutural e, portanto, ficaram enquadrados no nível 4.
Alguns textos com pior desempenho foram enquadrados no nível 1, enquanto o nível zero foi atribuído a pouquíssimos casos, já que ele revela a existência de um texto sem articulação e com comprometimento do sentido.
Observou -se que a grande maioria dos estudantes não distribuiu as ideias em parágrafos, talvez devido ao pequeno número de linhas disponibilizadas para a resposta da questão ou, quem sabe, pela suposição de que não seria necessária essa divisão por não se tratar de um texto no modelo de uma redação dissertativo-argumentativa, como solicitado nos vestibulares.
Outro aspecto observado na estruturação textual foi a divisão em dois itens, provavelmente motivados pelo encaminhamento do enunciado das questões, que apresentavam dois pontos a serem detalhados : a questão 1 solicitava que os estudantes dissertassem sobre as consequências do transporte motorizado ( a ) e apresentassem ações de intervenção por parte do poder público ( b ) ; a questão 2 solicitava que os estudantes dissertassem sobre as causas da violência urbana ( a ) e os fatores para evitá -la ( b ).
Essa competência pode ser considerada como a mais problemática no que diz respeito ao desempenho linguístico dos estudantes, porque são muitos os problemas observados, desvios acumulados durante toda a formação escolar e que não se resolvem com um estudo autodidata, como acontece com regras ortográficas ou morfossintáticas : sequência justaposta de ideias sem encaixamentos sintáticos ; redução drástica de estruturas subordinadas, ao lado do aumento na frequência de estruturas coordenadas e absolutas ; redução no uso de conectores para expressar relações lógicas essenciais à construção do texto, substituídas pela exigência de inferência por parte do interlocutor para suprir a sua ausência ; emprego equivocado de operadores que não estabelecem relações lógicas coerentes entre ideias do texto ; emprego inadequado do pronome relativo ( com omissão da preposição ou a utilização de pronome inadequado, como “ onde ” ) ; repetição exaustiva de termos sem a utilização de procedimentos mais sofisticados de substituição ( hiperonímias, hiponímias, nominalizações, expressões metafóricas ) ; frases fragmentadas que comprometem a estrutura lógico-gramatical ; frases formadas apenas por oração subordinada, sem oração principal.
Um importante aspecto a destacar é o baixíssimo desempenho de uma parte dos estudantes em relação à estrutura formal do texto produzido, o que é extremamente preocupante ao se levar em conta que são graduandos em fase final de formação.
São frequentes os casos de desvios de estruturação frasal, com uso inadequado ou ausência de conectivos entre parágrafos e entre frases.
Em uma parte dos textos, falta um mínimo de textualidade e de domínio do registro padrão da língua.
Na verdade, observam -se relações linguísticas quase agramaticais, como as estabelecidas pela sequência de gerúndios sem o apoio de um ponto de partida para a organização das informações gramaticais e semânticas, ou seja, sem uma oração principal.
Quanto à utilização dos mecanismos de referenciação, deve -se destacar a ocorrência, em uma boa parte dos textos, de repetições de palavras ou expressões sem a utilização de termos sinônimos ou pronomes, como seria adequado.
Quanto à utilização dos sinais de pontuação, observou -se uma grande precariedade nos textos analisados.
É muito frequente a ocorrência de parágrafos sem marca interna de pontuação para separar os períodos.
Vale observar que não foi penalizada a ausência de vírgula para destacar locuções ou adjuntos adverbiais de pequena extensão deslocados de posição na frase, por ser um uso opcional.
São os seguintes os tipos de problemas encontrados :  vírgula : utilização de vírgula para separar o sujeito e o predicado ; ocorrência de apenas uma das vírgulas para separar uma palavra, uma expressão ou uma oração encaixada ; uso de vírgula no lugar do ponto para separar ideias que constituem períodos distintos ; ausência de vírgula para separar elementos de uma enumeração ; ausência de vírgula para separar oração adjetiva explicativa ou utilização inadequada para separar oração adjetiva restritiva ;  ponto e vírgula : utilização do ponto e vírgula no lugar de vírgula ;  ponto final : ausência de ponto final para separar períodos.
Aspectos morfossintáticos e vocabulares : Da mesma forma que nos aspectos anteriores os textos foram classificados em níveis de zero a quatro.
O desempenho dos estudantes ficou concentrado, majoritariamente, nos níveis 3 e 2, devido à grande ocorrência de problemas de regência e concordância.
O nível 4 foi atribuído a um número menor de textos.
Alguns textos com pior desempenho foram enquadrados no nível 1, enquanto o nível zero foi atribuído a pouquíssimos casos, já que ele revela a existência de um texto sem o respeito às mínimas exigências morfossintáticas da norma-padrão e com comprometimento do sentido.
Os resultados são muito transparentes em relação aos aspectos mais problemáticos do desempenho dos estudantes.
O desvio mais frequente, em relação à regência, é a falta do sinal indicativo da crase – isso revela que o usuário não tem consciência de que, sob a forma do termo “ a ”, existe a presença de uma preposição “ a ”, exigida pela regência do termo anterior.
Embora em outros exames, como o Enem, a falta de crase seja penalizada em acentuação, nesta correção esse desvio foi considerado no âmbito dos aspectos morfossintáticos.
Outro problema relacionado à regência verbal e à nominal, encontrado frequentemente nas questões, foi a ausência de preposição antes de pronome relativo, processo generalizado na modalidade oral da língua, em situações de registro informal.
Apesar da possibilidade de que essa alteração de regência se generalize no padrão escrito da Língua Portuguesa, como já está ocorrendo até em textos jornalísticos, o não emprego da preposição foi penalizado neste processo de avaliação.
Outro desvio muito frequente diz respeito aos processos de concordância verbal e de concordância nominal.
Quanto à concordância de número, observou -se ausência de marca ( com sujeito anteposto ou posposto ) ou uso indevido ( uso inadequado da marca de plural comandado pelo núcleo plural da locução adjetiva, apesar de o substantivo que funciona como núcleo do sintagma nominal estar no singular ).
Uma ocorrência que se destacou foi a ausência de acento circunflexo na forma plural do presente do indicativo dos verbos “ ter ” e “ vir ”, que foi considerada como um desvio na concordância verbal e não na acentuação gráfica.
Quanto à concordância de gênero, vários casos foram observados, normalmente no âmbito de sintagmas nominais longos, em que o adjetivo está afastado do substantivo.
Deve -se destacar uma ocorrência não observada no ENADE/2013 : o aparecimento da marca de plural em verbos ou adjetivos relacionados a núcleos substantivos no singular, evidenciando um processo de hipercorreção.
Quanto à questão da colocação pronominal, foram poucos os casos observados.
Apesar de serem aspectos relacionados à oralidade, concluiu -se que, no registro escrito formal, a maioria dos estudantes já incorporou regras como a não introdução da frase por um pronome oblíquo e a próclise na presença de um termo atrator.
Não se adotou, entretanto, o padrão excessivamente formal descrito pelas gramáticas normativas em relação à posição do pronome oblíquo em locuções verbais, já que esse uso está muito distante da prática cotidiana, até em textos mais formais.
Quanto aos aspectos vocabulares, alguns tipos de inadequação foram observados : expressões da oralidade apareceram em algumas respostas, mas sem maior relevância do ponto de vista quantitativo ; seleção vocabular incompatível com o contexto, gerando situações de falta de inteligibilidade ; falta de domínio de vocabulário mais abstrato e de maior complexidade, essencial ao desenvolvimento do texto de base dissertativa.
Várias marcas de oralidade foram identificadas, embora não em alta frequência : o uso do pronome relativo “ onde ” como relativo universal, falta de artigo definido antes de substantivo, repetição de palavras por falta de vocabulário, reduções como “ tá ”, “ pra ”, “ pro ”, “ prum ”, expressões informais. 3.3.2 Componente de Conhecimento Específico Na parte da prova relativa às questões discursivas no Componente de Conhecimento Específico ( Tabela 3.17 ), observa -se que a média foi bem mais baixa do que para as questões discursivas do Componente de Formação Geral.
Enquanto no Componente de Formação Geral a média para estudantes de Física ( Bacharelado e Licenciatura ) de todo o Brasil foi 50,9, na parte de Conhecimento Específico a média inadequado da marca de plural comandado pelo núcleo plural da locução adjetiva, apesar de o substantivo que funciona como núcleo do sintagma nominal estar no singular ).
Uma ocorrência que se destacou foi a ausência de acento circunflexo na forma plural do presente do indicativo dos verbos “ ter ” e “ vir ”, que foi considerada como um desvio na concordância verbal e não na acentuação gráfica.
Quanto à concordância de gênero, vários casos foram observados, normalmente no âmbito de sintagmas nominais longos, em que o adjetivo está afastado do substantivo.
Deve -se destacar uma ocorrência não observada no ENADE/2013 : o aparecimento da marca de plural em verbos ou adjetivos relacionados a núcleos substantivos no singular, evidenciando um processo de hipercorreção.
Quanto à questão da colocação pronominal, foram poucos os casos observados.
Apesar de serem aspectos relacionados à oralidade, concluiu -se que, no registro escrito formal, a maioria dos estudantes já incorporou regras como a não introdução da frase por um pronome oblíquo e a próclise na presença de um termo atrator.
Não se adotou, entretanto, o padrão excessivamente formal descrito pelas gramáticas normativas em relação à posição do pronome oblíquo em locuções verbais, já que esse uso está muito distante da prática cotidiana, até em textos mais formais.
Quanto aos aspectos vocabulares, alguns tipos de inadequação foram observados : expressões da oralidade apareceram em algumas respostas, mas sem maior relevância do ponto de vista quantitativo ; seleção vocabular incompatível com o contexto, gerando situações de falta de inteligibilidade ; falta de domínio de vocabulário mais abstrato e de maior complexidade, essencial ao desenvolvimento do texto de base dissertativa.
Várias marcas de oralidade foram identificadas, embora não em alta frequência : o uso do pronome relativo “ onde ” como relativo universal, falta de artigo definido antes de substantivo, repetição de palavras por falta de vocabulário, reduções como “ tá ”, “ pra ”, “ pro ”, “ prum ”, expressões informais. 3.3.2 Componente de Conhecimento Específico Na parte da prova relativa às questões discursivas no Componente de Conhecimento Específico ( Tabela 3.17 ), observa -se que a média foi bem mais baixa do que para as questões discursivas do Componente de Formação Geral.
Enquanto no Componente de Formação Geral a média para estudantes de Física ( Bacharelado e Licenciatura ) de todo o Brasil foi 50,9, na parte de Conhecimento Específico a média O Gráfico 3.11 representa a distribuição das notas nas questões discursivas no Componente de Conhecimento Específico segundo a opção de Bacharelado/Licenciatura.
As modas destas distribuições ocorrem no primeiro intervalo, [ 0 ; 10 ], para a Licenciatura e no intervalo ( 30 ; 40 ] para o Bacharelado.
A análise de cada uma destas questões será feita a seguir. foi 36,0.
A maior média deste componente foi obtida pelos estudantes da região Sul ( 44,0 ), e a menor, pelos da região Norte ( 27,7 ).
Quanto à variabilidade das notas, o desvio padrão de todo o Brasil foi 24,6.
O maior desvio padrão foi encontrado na região Centro-Oeste ( 26,0 ), e o menor, na região Norte ( 21,3 ).
A maior nota máxima foi obtida nas regiões Nordeste, Sudeste e Sul ( 100,0 ), enquanto a nota máxima encontrada na região Norte foi 86,7 e na região Centro-Oeste foi 93,3.
Além disso, a nota mínima ( 0,0 ) foi obtida por pelo menos um estudante em todas as regiões do Brasil.
A mediana do Brasil como um todo foi 35,0, a mesma obtida na região Centro-Oeste.
Nas demais regiões a mediana foi : 26,7 na região Norte, 31,7 na região Nordeste, 43,3 na região Sudeste e 45,0 na região Sul.
O Gráfico 3.11 representa a distribuição das notas nas questões discursivas no Componente de Conhecimento Específico segundo a opção de Bacharelado/Licenciatura.
As modas destas distribuições ocorrem no primeiro intervalo, [ 0 ; 10 ], para a Licenciatura e no intervalo ( 30 ; 40 ] para o Bacharelado.
A análise de cada uma destas questões será feita a seguir. 3.3.2.1 Análise da Questão Discursiva 3 do Componente de Conhecimento Específico Na questão 3 ( Bacharelado e Licenciatura ), cujos resultados aferidos encontram -se descritos na Tabela 3.18, a média dos estudantes de todo o Brasil foi 45,8.
A menor média nessa questão foi obtida pelos estudantes da região Norte ( 34,4 ), enquanto a maior média foi obtida na região Sul ( 55,1 ).
Quanto à variabilidade das notas, o desvio padrão de todo o Brasil foi 36,8.
O maior desvio padrão foi obtido na região Centro-Oeste ( 38,0 ), enquanto o menor foi obtido na região Norte ( 34,5 ).
A nota máxima, 100,0 pontos, foi alcançada por, pelo menos, um estudante em todas as regiões do Brasil, assim como a nota mínima ( 0,0 ).
A mediana do Brasil como um todo foi 50,0, a mesma obtida nas regiões Sudeste e Sul, enquanto nas demais regiões, a mediana foi 35,0.
O Gráfico 3.12 mostra a distribuição das notas na questão discursiva 3 do Componente de Conhecimento Específico da área de Física para estudantes do Bacharelado e da Licenciatura.
Essa distribuição tem moda nas questões em branco para a Licenciatura e, no intervalo ( 90 ; 100 ], para o Bacharelado. 3.3.2.2 Comentários sobre as respostas à Questão Discursiva 3 A Questão, comum aos estudantes dos cursos de Licenciatura e Bacharelado, avaliava conteúdos de física básica, especificamente termodinâmica ( física térmica ), em nível adequado aos anos iniciais do ensino superior, sendo um tema abordado no ensino médio, essencialmente com a mesma profundidade.
O nível de dificuldade para a compreensão da situação-problema proposta, e de sua resolução, era simples, envolvendo noções conceituais sobre a dilatação térmica.
Exemplos similares ao proposto, com a utilização de barras metálicas e apresentação de situações análogas, estão presentes em grande número de livros didáticos e exames de ensino médio.
A solução da questão não envolvia qualquer cálculo, exigia apenas a compreensão conceitual do fenômeno da dilatação térmica e a aplicação dos conceitos relevantes à situação proposta.
Em relação às demais questões discursivas da parte de Conhecimentos Específicos, essa questão pode ser considerada fácil.
O padrão de respostas confirma a ideia de que a solução é simples.
A questão poderia ser resolvida em tempo adequado, e o espaço para apresentação da solução mostrou -se apropriado.
As perguntas poderiam ser respondidas de forma direta, porém os comandos deram origem a incompreensões por parte dos estudantes.
A afirmação do primeiro item – “ Explique uma forma de separar essas duas peças metálicas (... ) ” não garantiu que os estudantes compreendessem que deveriam propor uma forma de separar as duas peças metálicas e justificar sua proposta.
No segundo item, pedia -se a indicação de “ duas variáveis ( ou condições ) das quais o seu [ indicado pelo estudante no item ‘ a ’ ] método ” de separação poderia depender, o que possibilitou interpretações diversas, como será exposto na sequência.
O padrão de respostas indicava uma resposta discursiva, conceitual.
De modo geral, os estudantes indicaram o conceito correto para o encaminhamento da solução e o coeficiente de dilatação, mas a maioria não indicou o processo que permitiria a proposta de resolução da situação colocada.
Em geral, observou -se uma incompletude no raciocínio.
Em um conjunto significativo de respostas foi identificado o seguinte erro conceitual : supor que o processo permitiria o aquecimento ou resfriamento apenas de uma das peças, como se não houvesse condução térmica ou que esta fosse extremamente lenta.
Com relação à especificação da solução para o segundo item, verificou -se alguma incompreensão relativa ao que era esperado pelo padrão de respostas.
Quase Quase sempre os estudantes indicaram apenas a variável coeficiente de dilatação, apesar de terem sido solicitadas explicitamente, pelo menos, duas variáveis.
Os estudantes apresentaram desempenho razoável nesta questão, sendo o desempenho dos de Licenciatura mais fraco do que o dos de Bacharelado.
Tendo em vista o que era cobrado na questão, esperar -se-ia um desempenho melhor de ambos os grupos.
No grupo de respostas fracas ( notas de zero a 30 pontos ), verificou -se desconhecimento ou incompreensão do fenômeno da dilatação térmica e das variáveis que o produzem.
Surpreendeu encontrar entre estudantes em Física o total desconhecimento do assunto dilatação térmica e de conceitos simples associados à física térmica.
As respostas medianas ( notas de 35 a 70 pontos ), em geral, indicavam que era necessário aquecer ou esfriar o sistema, não explicando adequadamente como o procedimento deveria ser feito, ou indicavam apenas uma das variáveis que influem na dilatação.
As respostas consideradas boas ( notas iguais ou superiores a 75 pontos ), são respostas de estudantes que mostraram que conhecem o fenômeno de dilatação térmica, entendem como aplicá -lo a um caso concreto, e conhecem fatores que influem no fenômeno, mesmo cometendo algum pequeno erro.
A resposta mais frequentemente encontrada correspondia, no primeiro item, a informar que bastaria aquecer ( ou esfriar ) as peças, ou uma das partes da peça, sem discussão do processo envolvido ( distinguindo as possibilidades relativas ao maior ou menor coeficiente de dilatação na parte interna ou externa ).
Como já sinalizado, o principal erro cometido foi não considerar o processo em si.
No entanto, também se apresentou de forma importante a solução em que o estudante considera que o aquecimento ( ou esfriamento ) de uma única das partes não implicaria aquecimento ( ou esfriamento ) da outra, como se fosse possível não existir condução térmica.
No segundo item, a solução majoritária foi informar que a variável relevante era o coeficiente de dilatação térmica, sem uma segunda condição, mesmo o enunciado solicitando “ indique ao menos duas variáveis ”.
Esses foram os tipos de resposta encontrados, o que não configura uma grande diversidade de soluções.
Chama a atenção a dificuldade em argumentar de forma condicional ( por exemplo, foi difícil encontrar a argumentação"se o coeficiente de dilatação da parte externa for maior, então devemos aquecer esta parte"), algo 3.3.2.3 Análise da Questão Discursiva 4 do Componente de Conhecimento Específico A Tabela 3.19 contém as informações relativas à questão 4 do conjunto de questões do Componente de Conhecimento Específico, comum ao Bacharelado e à Licenciatura.
O desempenho dos estudantes de todo o Brasil nesta questão foi intermediário dentre as três questões discursivas desse componente.
A média geral do Brasil foi 37,4, sendo a menor média registrada na região Norte ( 30,1 ) e a maior, na região Sul ( 45,3 ).
A nota máxima ( 100,0 ) foi atingida por pelo menos um concluinte em todas as regiões do Brasil, assim como a nota mínima ( 0,0 ).
A mediana do Brasil, como um todo, foi 40,0, sendo maior para as regiões Sudeste, Sul e Centro-Oeste ( 50,0 em cada ) e menor para as regiões Norte e Nordeste ( 30,0 em ambas ).
O Gráfico 3.13 representa a distribuição de notas na questão discursiva 4, no Componente de Conhecimento Específico para estudantes do Bacharelado e da Licenciatura.
A maior frequência, tanto para estudantes de Bacharelado quanto para Licenciatura, ocorre na categoria “ em branco ”.
Observa -se que uma maior proporção de estudantes deixou essa questão em branco, comparando -se à questão discursiva 3.
Desconsiderando -se as respostas em branco a moda seria : para os estudantes do Bacharelado, o intervalo ( 60 ; 70 ], e, para os da Licenciatura, o intervalo ( 20 ; 30 ]. 3.3.2.4 Comentários sobre as respostas à Questão Discursiva 4 A questão foi proposta para os formandos dos cursos de Licenciatura e Bacharelado.
Referia -se a um tema pouco privilegiado em cursos de nível superior, focado em questões vinculadas às relações entre os efeitos das atividades científicas, tecnologias e a vida em sociedade.
No entanto, trata -se de um tema bastante relevante do ponto de vista do letramento científico da população, de acordo com a definição apresentada pelo Consórcio do PISA ( “ a capacidade de usar o conhecimento científico, de identificar questões e chegar a conclusões baseadas em evidências, de modo a compreender e a ajudar na tomada de decisões a respeito do mundo natural e as mudanças causadas a ele pela atividade humana ” ).
O que vem a confirmar que o conteúdo abordado era apropriado.
No entanto, uma consulta a fontes do Instituto de Radioproteção e Dosimetria da Comissão Nacional de Energia Nuclear - CNEN revela algumas imprecisões no texto base : dados coletados em diversos desses documentos não coincidem com os expostos no texto apresentado pelo enunciado da questão.
Inicialmente, a fonte de Césio do aparelho encontrado continha 28 g de cloreto de césio e 63 g de aglutinante, com uma massa total de 91 g.
Na abertura ( setembro de 1987 ) a atividade era de 1375 Cu, a massa de 91 g, e a massa de cloreto de césio ( CsCl ) valia 19,26 g.
Foram 90 \% liberados em Goiânia, portanto 17 g.
O parágrafo seguinte do referido texto, “ Por precaução (... ) ”, não é exato.
Só foi armazenado o que efetivamente estava contaminado.
Os materiais foram separados em dois depósitos, e não em dois containers.
O primeiro depósito foi utilizado para o material de baixa intensidade de radiação, e não de “ baixa concentração ”.
O segundo depósito foi usado para materiais de alta intensidade como os tambores com o material.
Além disso, não foi encontrada nenhuma referência que mencionasse uma camada de chumbo de 10 cm.
O primeiro item a ser respondido pelos estudantes trazia uma pergunta relativa à diferença entre contaminação e irradiação, dois conceitos importantes – o item está claro e adequado.
O segundo item mencionava uma equação, válida para núcleos.
No entanto, a vida média do Césio-137 é 43 anos e não 30 anos.
A meia-vida do Césio-137 é 30 anos, e os conceitos de vida média e meia-vida não são os mesmos.
A vida média é o tempo que a amostra leva para perder 1/e de seu número original, e a meia-vida é o tempo que leva a perder a metade ( as duas medidas estão relacionadas pela equação vida média = meia-vida / ln ( 2 ) ).
Mesmo com essas imprecisões, o item ‘ b ’ avaliava uma competência básica prevista nas Diretrizes Curriculares Nacionais para o curso : “ diagnosticar, formular e encaminhar a solução de problemas físicos, experimentais ou teóricos, práticos ou abstratos, fazendo uso dos instrumentos laboratoriais ou matemáticos apropriados ”.
E avaliava uma habilidade básica, geral, a de “ utilizar a matemática como uma linguagem para a expressão dos fenômenos naturais ” ; no caso particular, a questão exigia que o estudante fosse capaz de interpretar e fazer estimativas numéricas para resultados relevantes, a partir de equações matemáticas previstas pelos modelos físicos.
Mais até : este subitem exigia do estudante conectar a informação contida no texto com uma expressão matemática também fornecida no texto, e realizar um cálculo simples, uma estimativa, fazendo conexão entre as duas informações.
O terceiro item, letra ‘ c ’, apresentava uma formulação não inteiramente adequada, e foi anulada pela Comissão Assessora.
A questão estava adequada do ponto de vista do tempo de prova, e o espaço para apresentação da solução também foi apropriado.
Pode ser considerada como uma questão de dificuldade média.
O esperado pelo padrão de respostas adotado referia -se, no primeiro item, a uma diferenciação clara entre o conceito de contaminação ( no qual o indivíduo torna -se ele próprio um emissor ) e de irradiação.
Os estudantes demonstraram, em sua maioria, não ter clareza sobre a diferença entre os conceitos.
A formulação mais encontrada foi a de que contaminação corresponde a contato físico com um material contaminante, e que irradiação se confunde com radiação.
As respostas a esse item revelaram um letramento científico dos estudantes bastante pequeno, envolvendo um tópico de grande interesse para a vida, e que costuma ser veiculado em notícias ( como no caso do recente acidente em Fukushima ).
A diversidade das respostas não foi grande : respostas com a definição e diferenciação dos conceitos correta, respostas com o conceito de contaminação associado a contato físico ou à biologia apenas ( vírus, por exemplo ), respostas que consideravam irradiação como sinônimo de radiação.
O item ‘ b ’ exigia uma justificativa e uma argumentação com base em linguagem matemática, ou utilizando conceitos matematizados.
Os estudantes, em geral, conseguiram realizar os dois primeiros passos para a solução do item, sem completá -lo.
Poucos estudantes utilizaram a ideia do conceito de meia-vida para chegar à resposta de forma direta.
Em geral, conseguiram apresentar soluções que indicavam a compreensão da leitura ( o tempo decorrido entre o acidente e a data da prova era de quase 30 anos ), mas não desenvolveram a estimativa solicitada.
Cabe mencionar que no item ‘ b ’ não foi exigido um número exato, e sim a indicação de qualquer tipo de estimativa numérica que desse resultado de valor razoável ( inferior ao valor original e não nulo ).
Esse tipo de estimativa é corriqueiro em diversos tópicos do curso de Física, e foi surpreendente não ter sido desenvolvido pelos estudantes.
Devido à anulação do item ‘ c ’ dessa questão, que valia 30 pontos, todos os estudantes tiveram grau igual ou superior a 30.
Portanto, para adoção da classificação das notas em fracas, medianas e boas, nessa questão, é necessária uma adequação, lembrando -se de que a nota 30 corresponde a ter tirado zero nos itens ‘ a ’ e ‘ b ’.
Cerca de 20 \% dos estudantes tiveram nota 30 e, nesse caso, eles não informaram a definição de contaminação e irradiação, e também não conseguiram sequer calcular o tempo decorrido entre o ano do episódio de Goiânia e o ano atual.
As notas entre 35 e 70 correspondem aos estudantes que, em geral, apresentaram uma ligeira informação sobre os conceitos, mencionando, por exemplo, que contaminação ocorre quando há contato físico, ou avaliando corretamente o tempo decorrido entre o episódio e a data da prova, ou até mesmo indicando a expressão analítica para a solução, sem fazer a estimativa desejada.
Os graus iguais ou superiores a 75 correspondem a estudantes que demonstraram conhecer os conceitos de contaminação e irradiação, às vezes com pequenas dificuldades em sua apresentação ( por exemplo, mencionando irradiação como radiação e contaminação como contato físico ) e que não receberam a pontuação integral, ou mesmo os que não completaram a estimativa solicitada no item ‘ b ’.
É interessante mencionar também que o desempenho dos estudantes da Licenciatura revelou -se, novamente, inferior ao dos estudantes de Bacharelado.
A questão revelou uma deficiência na formação dos estudantes em relação a um conteúdo associado à capacidade de compreensão de conceitos importantes no que se refere à utilização de artefatos nucleares, para fins médicos.
A correção mostrou que os concluintes dos cursos de Física não dominam os dois conceitos, o que foi inesperado, tendo em vista o destaque dado a episódios deste tipo ( como por exemplo, a longa discussão recente sobre o acidente de Fukushima ).
Outra deficiência preocupante foi a dificuldade revelada de obtenção de valores numéricos estimados a partir de uma expressão analítica bastante comum em diferentes situações que descrevem fenômenos físicos.
O resultado do desempenho na questão 4, portanto, evidencia a necessidade de reflexão das coordenações de curso sobre a formação que vem sendo recebida pelos estudantes dos cursos de Física.
Apesar de não abordar um tema de presença dominante no currículo, a existência dessa questão na prova acaba por provocar reflexões interessantes sobre a formação dos estudantes.
É importante frisar que esse tipo de questão é fundamental para a compreensão das competências e conceitos que deveriam ser construídos pelos estudantes em sua formação.
No entanto, preocupa a imprecisão do texto apresentado, que fez com que o item ‘ c ’ precisasse ser anulado, o que prejudicou a avaliação por meio dos resultados. 3.3.2.5 Análise da Questão Discursiva 5 do Componente de Conhecimento Específico ( Bacharelado ) A Tabela 3.20 contém as informações relativas à questão 5 do conjunto do Componente de Conhecimento Específico ( Bacharelado ).
O desempenho dos estudantes nessa questão foi inferior ao das questões 3 e 4.
A nota média dos estudantes de todo o Brasil foi 25,8.
A maior média foi registrada na região Centro-Oeste ( 33,1 ), enquanto a menor média foi registrada na região Norte ( 20,4 ).
Quanto à variabilidade das notas, o desvio padrão dos estudantes do Brasil, como um todo, foi 32,2.
Enquanto o maior desvio foi encontrado na região Centro-Oeste ( 33,4 ), o menor foi encontrado na região Norte ( 29,1 ).
A nota mínima ( 0,0 ) foi a mesma em todas as regiões.
A nota máxima 100,0 foi atingida, por pelo menos um aluno, nas regiões Nordeste, Sudeste e Centro-Oeste.
Na região Norte a nota máxima foi 80,0 e na região Sul foi 95,0.
Para o conjunto de estudantes de Bacharelado em Física do Brasil, a mediana foi 10,0, a mesma obtida nas regiões Sudeste e Sul.
Na região Centro-Oeste a mediana foi maior ( 25,0 ), enquanto nas regiões Norte ( 0,0 ) e Nordeste ( 5,0 ) a mediana foi menor.
O Gráfico 3.14 representa a distribuição das notas da questão discursiva 5 no Componente de Conhecimento Específico ( Bacharelado ).
Mais uma vez, destaca -se o grande número de estudantes que deixaram a questão em branco, a categoria modal. 3.3.2.6 Comentários sobre as respostas à Questão Discursiva 5 ( Bacharelado ) A questão tratava da \textbf{observação} de dados experimentais relativos a medidas feitas em um circuito elétrico simples.
No que diz respeito às Diretrizes Curriculares Nacionais, a questão é inteiramente adequada.
O nível de profundidade é razoável, exigindo do estudante de Bacharelado em Física uma distinção não óbvia e uma síntese não usual, a partir de dados de resultados experimentais.
O enunciado da questão era claro.
Havia uma pequena dificuldade para a interpretação da figura, já que a curva iii foi colocada muito próxima tanto à curva i quanto à curva ii.
Também, ao contrário do que é usual no ensino médio e no ciclo básico do ensino superior, o gráfico apresentado foi de I vs V em vez do gráfico V vs I, o que levou, desnecessariamente, a algumas interpretações incorretas.
A questão estava adequada do ponto de vista do tempo de prova, e o espaço para apresentação da solução também foi apropriado.
Pode ser considerada como uma questão de dificuldade média.
No que se refere às respostas dadas pelos estudantes, observou -se que a generalidade da pergunta –"discuta a natureza física dos materiais [... ], dê exemplos [... ] e cite aplicações [... ]"– trouxe dificuldades, por falta de delimitação objetiva do que se esperava como resposta.
O padrão de respostas, porém, era abrangente o bastante para possibilitar a análise das respostas dadas pelos estudantes para a questão.
As respostas foram muito diversas, em função da generalidade comentada.
A questão também permitia respostas muito amplas, a partir da interpretação de um gráfico ( apresentado de forma não inteiramente usual ).
Nem todos os estudantes foram capazes de apresentar uma classificação da natureza física dos materiais da forma indicada no padrão de respostas.
Em geral, mencionavam resistores ôhmicos e não ôhmicos, não fornecendo a sequência condutor / semi-condutor / isolante.
Os exemplos sempre foram incompletos ou não apresentados, e as aplicações, muitas vezes, foram poucas.
A linguagem adotada pelos estudantes foi adequada.
Na correção da questão não se identificou um padrão para estas respostas fracas ( de zero a 30 pontos ).
Os erros foram bastante variados.
As notas medianas ( notas entre 35 e 70 ), em geral, correspondiam a caracterizações parciais dos materiais ( nem todos corretos ), ou a caracterizações sem exemplos nem aplicações.
A questão 5 do curso de Bacharelado era adequada à formação dos estudantes, e revelou a dificuldade que estes encontram em fazer inferências a partir de informações fornecidas em forma de gráfico e relacioná -las com os conteúdos, em geral, apresentados de forma teórica.
A formulação era bastante aberta, e isso permitiu múltiplas interpretações, o que estava contemplado no padrão de resposta.
No entanto, se a questão fosse menos aberta, a avaliação da aprendizagem seria também mais precisa. 3.3.2.7 Análise da Questão Discursiva 5 do Componente de Conhecimento Específico ( Licenciatura ) A Tabela 3.21 contém as informações relativas à questão 5 do conjunto do Componente de Conhecimento Específico ( Licenciatura ).
O desempenho dos estudantes nessa questão foi inferior ao das questões 3 e 4.
A nota média dos estudantes de todo o Brasil foi 24,6.
A maior média foi registrada na região Sul ( 34,4 ), A nota mínima ( 0,0 ) e a nota máxima ( 100,0 ) foram as mesmas em todas as regiões.
Para o conjunto de estudantes de Licenciatura em Física do Brasil, a mediana foi 0,0, o que também ocorreu nas regiões Norte, Nordeste e Centro-Oeste.
Ou seja, pelo menos a metade dos estudantes de todo o Brasil tirou nota zero nessa questão.
Apenas nas regiões Sudeste e Sul a mediana foi diferente de zero ( 20,0 ).
O Gráfico 3.15 representa a distribuição das notas da questão discursiva 5 no Componente de Conhecimento Específico ( Licenciatura ).
Mais uma vez, destaca -se o grande número de estudantes que deixaram a questão em branco ( categoria modal ), seguido de perto por aqueles que tiraram nota zero para a resposta apresentada. 3.3.2.8 Comentários sobre as respostas à Questão Discursiva 5 ( Licenciatura ) A questão apresentada aos estudantes de Licenciatura solicitava a elaboração de um “ texto dissertativo contendo uma crítica sobre a forma de utilização do experimento adotado pelo professor e uma sugestão de alternativa de procedimento ” [ grifo nosso ].
Portanto, o estudante deveria ler e analisar um texto que compunha o enunciado da questão para apresentar sua reflexão crítica.
O referido texto descrevia um experimento pouco adequado ao fim que se propunha, apresentando um experimento com muitas repetições, com dificuldades técnicas para sua realização e inadequado ao objetivo de aprendizagem.
A questão era clara, a linguagem utilizada era adequada, mas o conteúdo abordado inadequado em relação às Diretrizes Curriculares Nacionais.
A solicitação feita seria mais adequada a um curso de pós-graduação em ensino de Física, no qual as leituras teóricas sobre o papel e os formatos de laboratórios são exigência fundamental.
No entanto, a formação inicial do professor, em geral, não aborda essas leituras e, certamente, o fato de no experimento relatado o professor utilizar “ roteiros fechados ” não é o principal problema da situação apresentada.
Há problemas na modelagem da situação experimental, na realização do experimento, na proposta de repetições, entre outros.
Os conhecimentos exigidos pela questão estão desalinhados quanto ao nível de profundidade esperado pelo padrão de respostas.
Houve muitas interpretações diferentes do enunciado, com diversidade de respostas e poucas soluções ligadas ao padrão adotado.
A questão pode ser considerada muito difícil para os estudantes de Licenciatura.
A exigência prevista no padrão de respostas revelou -se inatingível pelos estudantes.
De fato, a crítica aos roteiros fechados prevista no padrão de respostas é assunto não dominado pelos estudantes, e isso não era inesperado, tendo em vista que a literatura que apresenta esta discussão não é usualmente apresentada em cursos de graduação.
Mais do que isso, a formação nas escolas do país é, em grande parte, realizada sem ajuda de laboratórios, e por isso exigia -se que o estudante superasse uma primeira dificuldade : considerar de fato importante ensinar física com atividades de laboratório e experimentais.
Superada essa dificuldade, intrínseca à sua própria formação, deveria ser capaz de fazer uma crítica ao uso proposto na situação apresentada do laboratório.
Destaca -se, ainda, que a experiência proposta no texto apresenta dificuldades conceituais ( a construção do modelo ) e de procedimento ( materiais e métodos ).
A diversidade de respostas encontrada foi muito alta, em função das características do próprio enunciado.
Poucos estudantes mencionaram críticas ao formato dos roteiros ; foram feitas muitas críticas e sugestões de mudanças, principalmente sobre os métodos, os materiais, o uso ( ou a falta de uso ) de laboratórios entre outros.
Muitas das respostas revelaram boa capacidade de redação, mesmo não apresentando exatamente o que era sugerido pelo padrão de respostas.
O desempenho dos estudantes foi muito fraco e ficou patente que a quase totalidade dos estudantes não entendeu o que se esperava deles.
Como na questão 5 do curso de Bacharelado, houve muita diversidade nas soluções apresentadas para a questão 5 do curso de Licenciatura.
Como já comentado, a questão possibilitou diversas interpretações, mas não, majoritariamente, a esperada no padrão de respostas.
As respostas referiam -se a melhorias no método de coleta de dados, à modelagem do experimento, a mudanças nos materiais, todas continham críticas ao experimento de tipos diversos.
A questão 5 do curso de Licenciatura era muito aberta, e o padrão de resposta se mostrou inadequado, por causa da própria formulação da questão.
De fato, não seria possível prever todas as possibilidades de resposta a serem consideradas corretas ou parcialmente corretas.
A questão se revelou inadequada para avaliar a formação de Licenciados no país, não porque o tema seja pouco importante ou pouco adequado, mas porque ele pressupõe uma maturidade no uso de laboratórios e experimentos em sala de aula à qual a realidade das escolas brasileiras, incluindo as universidades, não está adequada.
Para finalizar, a correção da parte discursiva da prova do ENADE/2014 de Física possibilitou a avaliação da adequação das Diretrizes Curriculares Nacionais e a formação dos estudantes, revelando algumas dificuldades e inadequações, particularmente nos cursos de Licenciatura, que exigem reflexões.
A formação de profissionais na área de Física, tanto pesquisadores físicos quanto a formação de professores para a educação básica apresenta -se não satisfatória, no global, tendo em vista o observado nesta prova.
Evidentemente, há estudantes que revelam em suas soluções uma formação boa e adequada ; no entanto, a maior parte dos concluintes demonstrou ter sérias dificuldades com a argumentação, com a escrita de suas ideias e com a compreensão de conceitos simples, a maior parte deles com relação direta com a utilização na vida cotidiana. 3.3.3 Considerações Finais As Diretrizes Curriculares da Área de Física pressupõem, para o perfil do egresso do curso, “ um profissional que, apoiado em conhecimentos sólidos e atualizados em Física, deve ser capaz de abordar e tratar problemas novos e tradicionais e deve estar sempre preocupado em buscar novas formas do saber e do fazer científico ou tecnológico ” ( Parecer CNE/CES nº 1304/2001 ).
São também definidos quatro perfis específicos : o Físico-pesquisador, o Físico-educador, o Físico-tecnólogo, e o Físico-interdisciplinar.
Os cursos de graduação em Física, tanto diplomando bacharéis quanto licenciados, devem ter em mente a formação para um desses perfis.
As competências exigidas para esses profissionais envolvem um conjunto grande de competências e habilidades, algumas das quais não podem ser avaliadas por meio de provas escritas, tanto dissertativas quanto objetivas ( como por exemplo, dentre as competências essenciais, desenvolver uma ética de atuação profissional, ou, dentre as habilidades, utilizar os diversos recursos de informática, dispondo de noções de linguagem computacional, apresentar resultados científicos em distintas formas de expressão, ou mesmo conhecer e absorver novas técnicas, métodos ou uso de instrumentos ).
Portanto, as provas aplicadas avaliam algumas das competências e habilidades previstas para os perfis de egressos constantes das Diretrizes Nacionais Curriculares.
Sem fazer alguma avaliação das questões objetivas, tanto no que se refere às suas qualidades ou defeitos quanto no que se refere à relação com o que estava sendo avaliado pela prova do ENADE/2014, busca -se, a seguir, compreender a relação de tais questões com as Diretrizes Curriculares Nacionais do curso. dentre as habilidades, utilizar os diversos recursos de informática, dispondo de noções de linguagem computacional, apresentar resultados científicos em distintas formas de expressão, ou mesmo conhecer e absorver novas técnicas, métodos ou uso de instrumentos ).
Portanto, as provas aplicadas avaliam algumas das competências e habilidades previstas para os perfis de egressos constantes das Diretrizes Nacionais Curriculares.
Sem fazer alguma avaliação das questões objetivas, tanto no que se refere às suas qualidades ou defeitos quanto no que se refere à relação com o que estava sendo avaliado pela prova do ENADE/2014, busca -se, a seguir, compreender a relação de tais questões com as Diretrizes Curriculares Nacionais do curso.
A prova para os estudantes dos cursos de Bacharelado e Licenciatura em Física apresenta um conjunto de questões objetivas comuns ( de números 9 a 25 ) relativas a conteúdos específicos de Física, envolvendo algumas das competências previstas nas Diretrizes Curriculares Nacionais do curso ; especificamente, dominar princípios gerais e fundamentos da Física, estando familiarizado com suas áreas clássicas e modernas, e descrever e explicar fenômenos naturais, processos e equipamentos tecnológicos em termos de conceitos, teorias e princípios físicos gerais.
A prova para os estudantes do curso de Bacharelado em Física apresenta um conjunto adicional de dez questões objetivas ( de número 26 a 35 ), relativas a conteúdos específicos de Física, envolvendo tópicos melhor adequados a perfis de formação mais ligados a temas mais avançados de física, avaliando as bases para a utilização das habilidades previstas nas Diretrizes Curriculares Nacionais : utilizar a matemática como uma linguagem para a expressão dos fenômenos naturais, propor, elaborar e utilizar modelos físicos, reconhecendo seus domínios de validade, e utilizar a linguagem científica na expressão de conceitos físicos, na descrição de procedimentos de trabalhos científicos e na divulgação de seus resultados.
A prova para os estudantes do curso de Licenciatura em Física apresenta um conjunto adicional de 10 questões objetivas ( de números 26 a 35 ) relativas a conteúdos vinculados ao trabalho docente, como o conhecimento dos Parâmetros Curriculares Nacionais para o Ensino Médio, as metas do Plano Nacional de Educação, aspectos relativos à elaboração de um currículo com ênfases mais atuais ( por exemplo, com uso de Tecnologias de Informação e Comunicação, história e filosofia da ciência e outros ) para a atuação profissional na educação formal básica, possibilitando avaliação das bases conceituais de habilidades específicas para a formação do professor especificadas nas Diretrizes Curriculares Nacionais : o planejamento e o desenvolvimento de diferentes experiências didáticas em Física, reconhecendo os elementos relevantes às estratégias adequadas e a elaboração ou adaptação de materiais didáticos de diferentes naturezas, identificando seus objetivos formativos, de aprendizagem e educacionais.
A parte discursiva da prova foi composta de três questões, duas comuns aos estudantes de Bacharelado e Licenciatura ( as questões discursivas 3 e 4 ) e uma específica ( questão discursiva 5 ) para a prova de estudantes de Bacharelado e para estudantes de Licenciatura.

% matched lemmas: observação
\end{textsample}
